\documentclass[11pt]{article}
\usepackage[T1]{fontenc}
\usepackage[utf8]{inputenc}
\usepackage{lmodern}
\usepackage[margin=1in,headheight=15pt]{geometry}
\usepackage{amsmath,amssymb,amsthm,mathtools,mathrsfs}
\usepackage{microtype,booktabs,array,longtable,enumitem}
\usepackage{xcolor}
\usepackage[colorlinks=true,linkcolor=blue!40!black,citecolor=green!30!black,urlcolor=blue!40!black]{hyperref}
\usepackage[nameinlink,noabbrev]{cleveref}
\usepackage{fancyhdr}
\pagestyle{fancy}
\fancyhf{}
\fancyhead[L]{\small Trigonometry on the surcomplex plane}
\fancyhead[R]{\small\thepage}
\renewcommand{\headrulewidth}{0.3pt}
\setlength{\parskip}{0.25em}
\setlength{\emergencystretch}{2em}
\setlist[enumerate]{itemsep=0.2em,topsep=0.3em}
\setcounter{tocdepth}{2}
\numberwithin{equation}{section}
\newtheorem{theorem}{Theorem}[section]
\newtheorem{lemma}[theorem]{Lemma}
\newtheorem{proposition}[theorem]{Proposition}
\newtheorem{corollary}[theorem]{Corollary}
\theoremstyle{definition}
\newtheorem{definition}[theorem]{Definition}
\newtheorem{example}[theorem]{Example}
\newtheorem{remark}[theorem]{Remark}
\newcommand{\No}{\mathbf{No}}
\newcommand{\SC}{\mathbf{No}[i]}
\newcommand{\R}{\mathbb R}
\newcommand{\C}{\mathbb C}
\newcommand{\Q}{\mathbb Q}
\newcommand{\Z}{\mathbb Z}
\newcommand{\N}{\mathbb N}
\newcommand{\OO}{\mathcal O}
\newcommand{\mm}{\mathfrak m}
\newcommand{\HH}{\mathcal H}
\newcommand{\TT}{\mathbb T}
\newcommand{\eps}{\varepsilon}
\newcommand{\st}{\operatorname{st}}
\newcommand{\supp}{\operatorname{supp}}
\newcommand{\re}{\operatorname{Re}}
\newcommand{\im}{\operatorname{Im}}
\newcommand{\cis}{\operatorname{cis}}
\newcommand{\Arg}{\operatorname{Arg}}
\newcommand{\arccot}{\operatorname{arccot}}
\newcommand{\arsinh}{\operatorname{arsinh}}
\newcommand{\arcosh}{\operatorname{arcosh}}
\newcommand{\artanh}{\operatorname{artanh}}
\newcommand{\Ex}{\operatorname{Exp}_{\mm}}
\newcommand{\Logm}{\operatorname{Log}_{\mm}}
\newcommand{\Sin}{\operatorname{Sin}}
\newcommand{\Cos}{\operatorname{Cos}}
\newcommand{\Tan}{\operatorname{Tan}}
\newcommand{\fin}{\operatorname{fin}}
\newcommand{\sh}{^{\sharp}}
\newcommand{\HInt}{\mathop{\overset{\scriptscriptstyle H}{\int}}\nolimits}
\newcommand{\abs}[1]{\left|#1\right|}
\newcommand{\ang}[2]{\angle(#1,#2)}
\newcolumntype{L}[1]{>{\raggedright\arraybackslash}p{#1}}
\hypersetup{pdftitle={Trigonometry on the Surcomplex Plane},pdfauthor={},pdfsubject={Canonical finite angles, non-Archimedean triangle geometry, trigonometric polynomials, and infinite phases}}
\begin{document}
\begin{titlepage}
\centering
\vspace*{1.1cm}
{\LARGE\bfseries Trigonometry on the\\[0.2em] Surcomplex Plane\par}
\vspace{0.65cm}
{\Large Canonical Angles, Infinite Triangles,\\[0.15em]
Infinitesimal Contact, and Hahn-Analytic Oscillation\par}
\vspace{0.8cm}
{\large September 21, 2026\par}
\vspace{0.65cm}
\begin{minipage}{0.94\textwidth}
\begin{abstract}
The surcomplex plane is the algebraically closed class field $\SC$, with Euclidean norm $|x+iy|=\sqrt{x^2+y^2}$. Its lengths can be infinite or infinitesimal, but every direction has a finite surreal angle. We make this observation into a trigonometric theory. Hahn summation gives canonical sine and cosine on finite surreal arguments and identifies the full unit circle with the group of finite angles modulo $2\pi\Z$. The circle splits into an ordinary phase and an additive infinitesimal phase; its angular and chordal distances have the same valuation.

We prove exact analogues of the angle-sum, sine, cosine, tangent, half-angle, Heron, Euler, angle-bisector, trigonometric Ceva, and Ptolemy theorems for triangles and circles at arbitrary surreal scales. Valuation formulas measure thinness, and nearly tangent intersections exhibit exact square-root splitting and a quantified loss of angular precision. A canonical complex-variable extension exists on the strip of surcomplex numbers with finite real part, even when the imaginary part is infinite. On this strip, trigonometric polynomials reduce to ordinary-degree algebraic polynomials; this yields sharp root counts, support-controlled cluster lifting, finite Fourier identities, and a Fej\'er--Riesz factorization over the surreal field. Hyperbolic sine and cosine laws hold in the full surreal Poincar\'e disk.

Finally, we classify the freedom in extending circular trigonometry to infinite real arguments. Such extensions exist and may be fine-locally analytic, but the usual local identities do not determine their infinite phases, and they necessarily acquire infinite periods. All infinite evaluations in the article have set-sized, locally finite Hahn support. Neither sequential convergence nor an unqualified transfer of real analysis is assumed.
\end{abstract}
\end{minipage}
\vfill
\begin{minipage}{0.94\textwidth}\small
\textbf{Status and provenance.} A self-contained development built on the summability conventions of the three supplied manuscripts. Classical identities and established surreal analytic machinery are identified as such. No priority claim or claim to solve a named published open problem is made. Exact symbolic checks accompany the proofs; they are not a formal verification of the theory.
\end{minipage}
\end{titlepage}
\begingroup
\small
\setlength{\parskip}{0pt}
\tableofcontents
\endgroup
\clearpage

\section{The organizing distinction: lengths versus angles}\label{sec:intro}
Let $z=x+iy$ with $x,y\in\No$ and $i^2=-1$. This is the meaning of a surcomplex number used throughout. The field structure and real closedness of $\No$ are established parts of surreal number theory \cite{Conway,Gonshor,Alling}. They imply that the elementary algebra of Euclidean geometry remains available: distances have nonnegative square roots, two-dimensional determinants measure oriented area, and finite systems of linear equations can be solved exactly.

The analytic difficulty is different. The expression
\[
 \sum_{n\ge0}\frac{(-1)^nx^{2n+1}}{(2n+1)!}
\]
cannot simply be declared to define $\sin x$ for every surreal $x$. At a nonzero ordinary real argument, infinitely many summands already have the same Hahn exponent. At a positive infinite argument, the monomial supports may run in the wrong direction. Ordinary real convergence does not repair either defect in the full surreal fine topology.

There is, nevertheless, a canonical and sufficient theory of geometric angles. If $z\ne0$, its normalized direction $z/|z|$ has finite coordinates. Such a direction has an ordinary standard part on the usual unit circle and an infinitesimal angular correction. Consequently, trigonometry of \emph{all} surcomplex triangles does not require a choice of $\sin\omega$.

\subsection{What comes from the supplied manuscripts}
The supplied analysis manuscript \cite{SourceAnalysis} provides the notation $t^\gamma=\omega^{-\gamma}$, strong summability, standard parts, Taylor evaluation on infinitesimal thickenings, coefficientwise contour functionals, and examples of nonunique global fine-analytic exponentials. We retain those conventions and reprove the elementary parts needed here.

The two archives \cite{SourceFinite,SourceIntersections} concern finite \emph{zero algebras} and analytic intersections, rather than Euclidean triangle geometry. Their relevant themes are conservation of multiplicity, positive-support coefficient recursion, and loss of precision near a singular Jacobian. Sections~\ref{sec:contact} and~\ref{sec:clusters} develop trigonometric instances of those themes. We do not assume their several-variable Noetherianity, division, or residue theorems to prove the geometric results below.

\subsection{Established background and the scope of the extension}
Restricted analytic evaluation on surreal fields is not new. Van den Dries--Ehrlich construct the restricted analytic and exponential expansion of $\No$ and prove an elementary-extension theorem \cite[Theorem 2.1]{vdDE}. Hahn-field support and algebraic-closure facts have established treatments, including Poonen \cite{Poonen}. Classical trigonometric identities are recorded, for example, in the NIST Digital Library of Mathematical Functions \cite{DLMF}. Their extensions below are proved directly, with domains stated explicitly.

The main additions of this article are the organization of these ingredients into an all-scale geometry, the exact angle/valuation dictionary, the elementary analysis of infinitesimal contact and trigonometric root clusters, and a clear boundary between canonical finite-angle trigonometry and freely chosen infinite phases. A proof of an analogue is not, by itself, evidence that its exact formulation is absent from all previous literature.

\begin{center}
\small
\begin{tabular}{L{0.23\textwidth}L{0.32\textwidth}L{0.35\textwidth}}
\toprule
Object & Domain used here & Governing construction\\
\midrule
Lengths and coordinates & All of $\No$ and $\SC$ & Ordered-field algebra and square roots\\
Circular angles & Finite $\theta\in\No$, modulo $2\pi\Z$ & Ordinary value plus infinitesimal Taylor series\\
Inverse tangent & Every $x\in\No$ & Direction of $1+ix$\\
Complex sine and cosine & $\{x+iy:x\text{ finite},\ y\in\No\}$ & Finite phase and Gonshor exponential\\
Trigonometric polynomials & Ordinary finite degree & Laurent-polynomial algebra on the unit circle\\
Sine at infinite real inputs & Only after an extension is specified & An additional additive-group character\\
\bottomrule
\end{tabular}
\end{center}

\section{Surreal scalars and support-controlled evaluation}\label{sec:foundations}
\subsection{Size, norm, valuation, and standard part}
All supports, polynomial degrees, coefficient families, and recursions in this article are sets. The notation $\No$ denotes a proper class; a class theory such as NBG may be used. Arguments involving any prescribed set of coefficients take place in a set-sized Hahn workspace. In particular, no sum has a proper-class index set.

Conjugation, norm, inner product, and oriented determinant are
\begin{align}
 \overline{x+iy}&=x-iy,& |z|&=\sqrt{z\bar z},\\
 \langle z,w\rangle&=\re(\bar zw),&[z,w]&=\im(\bar zw).
 \label{eq:dotcross}
\end{align}
Expanding coordinates proves
\begin{equation}\label{eq:gram}
 \langle z,w\rangle^2+[z,w]^2=|z|^2|w|^2.
\end{equation}
It follows that $|zw|=|z||w|$, $|\langle z,w\rangle|\le|z||w|$, and $|z+w|\le|z|+|w|$. Equality in the last inequality for nonzero $z,w$ occurs exactly when $w/z$ is positive real in $\No$. Here and below, ``real'' as a coordinate condition means belonging to $\No$; ``ordinary real'' means belonging to $\R$.

Fix a set-sized divisible ordered additive subgroup $\Gamma\subset\No$ containing the exponents under consideration. Write
\[
 F_\Gamma=\R((t^\Gamma)),\qquad
 K_\Gamma=\C((t^\Gamma)),\qquad t^\gamma=\omega^{-\gamma}.
\]
These embed canonically in $\No$ and $\SC$. The real field $F_\Gamma$ is real closed and $K_\Gamma$ is algebraically closed; alternatively one may take algebraic roots directly in the full class fields \cite{vdDE,Poonen}. For a nonzero normal form $a=\sum_{\gamma\in S}a_\gamma t^\gamma$, put
\[
 v(a)=\min S,\qquad v(0)=+\infty.
\]
Then $v(ab)=v(a)+v(b)$, $v(a+b)\ge\min(v(a),v(b))$, and $v(|a|)=v(a)$.

Let $\OO_{\R}$ be the finite surreals and $\mm_{\R}$ the infinitesimals:
\[
 \OO_{\R}=\{x:|x|<n\text{ for some }n\in\N_{>0}\},\quad
 \mm_{\R}=\{x:|x|<1/n\text{ for every }n\in\N_{>0}\}.
\]
Use subscripts $\C$ for the analogous subclasses of $\SC$. Every finite element has a unique decomposition
\begin{equation}\label{eq:standard}
 z=c+\eta,\qquad c=\st(z)\in\C,\quad \eta\in\mm_{\C}.
\end{equation}
On real elements $c\in\R$. Standard part is a ring homomorphism. An element of $\OO_{\R}$ is called \emph{appreciable} when its standard part is nonzero. Thus finite nonzero elements have valuation zero exactly when they are appreciable.

We write $a\prec b$ when $a/b$ is infinitesimal, for $b\ne0$. For a specified nonzero infinitesimal $h$, the notation $O(h^m)$ means a remainder whose quotient by $h^m$ is finite. It is an algebraic size statement, not an assertion of convergence as an ordinary integer tends to infinity.

\subsection{The support lemma}
\begin{definition}[Strong summability]
A set-indexed Hahn family $(A_j)_{j\in J}$ is strongly summable when the union of its supports is well ordered and every exponent occurs in the supports of only finitely many family members. Its sum is formed by finite addition at each exponent.
\end{definition}

\begin{lemma}[Hahn--Neumann support calculus]\label{lem:support}
If $S,T$ are well-ordered subsets of an ordered abelian group, then $S+T$ is well ordered and each exponent has finitely many representations $s+t$. If $T\subset\Gamma_{>0}$ is well ordered, the monoid
\[
 T^*=\{\tau_1+\cdots+\tau_k:k\ge0,\ \tau_j\in T\}
\]
is well ordered, and each exponent has only finitely many representations as a finite word in $T$, including the possible word lengths. In particular, for finitely many infinitesimals $\eta_1,\ldots,\eta_d$ and arbitrary ordinary complex coefficients $c_\alpha$, the family
\[
 \sum_{\alpha\in\N^d}c_\alpha\eta_1^{\alpha_1}\cdots\eta_d^{\alpha_d}
\]
is strongly summable. Products and formal substitutions with zero constant term may be regrouped coefficientwise.
\end{lemma}
\begin{proof}
The first assertion follows by extracting a nondecreasing subsequence in one coordinate: either a decreasing sequence of sums or infinitely many representations of one sum would then force the other coordinate to decrease. The positive-monoid assertion is the Neumann lemma \cite{Neumann}. One convenient proof uses the well-quasi-ordering of finite words in a well-ordered alphabet under subsequence embedding with coordinatewise increase. Such an embedding cannot decrease a sum of positive letters, and equality forces identical letters and identical length. This excludes both descending sums and infinitely many distinct nondecreasing words of a fixed sum; each word has finitely many permutations. Applying this to the union of the positive supports of the $\eta_j$ proves the last assertions. The well-quasi-ordering statement is the standard finite-word lemma used in this formulation of Neumann's support argument.
\end{proof}

The lemma applies even when the value group has arbitrary rank. For example, with $0<\lambda\ll\mu$ one can have $n\lambda<\mu$ for every ordinary $n$, while $\sum_{n\ge0}t^{n\lambda}$ is a perfectly valid Hahn series. Its partial sums need not converge in the valuation topology.

\subsection{Lifting ordinary analytic germs}
For an ordinary real-analytic function $f$ near $r\in\R$, define
\begin{equation}\label{eq:lift}
 f\sh(r+\eps)=\sum_{n\ge0}\frac{f^{(n)}(r)}{n!}\eps^n,
 \qquad \eps\in\mm_{\R}.
\end{equation}
There is an identical definition for ordinary holomorphic germs and complex infinitesimals. Formula~\eqref{eq:lift} means a Hahn sum; the ordinary value $f(r)$ is taken first. Formal identities among the germs lift to exact identities, because both sides have identical Taylor coefficients and \cref{lem:support} permits the substitutions.

\begin{proposition}[Local calculus and one-variable sign lifting]\label{prop:lift}
Taylor lifting preserves sums, products, composition where defined, and differentiation. It is the unique extension specified by the Taylor rule \eqref{eq:lift}. If an ordinary analytic $g$ is nonnegative on an ordinary closed interval $[a,b]$ and analytic on a neighborhood of it, then $g\sh(x)\ge0$ for every surreal $x\in[a,b]$.
\end{proposition}
\begin{proof}
The algebraic and compositional statements are formal Taylor identities justified by the support lemma. Re-expansion at $x$ gives
\[
 f\sh(x+h)=f\sh(x)+(f')\sh(x)h+h^2R_x(h)
\]
for infinitesimal $h$, with $R_x(h)$ finite on a sufficiently small neighborhood. This follows either from the convergent ordinary germ before substitution or directly from its Taylor coefficients and their standard parts. The difference quotient therefore tends to $(f')\sh(x)$ in the fine $\eps$--$\delta$ sense: for any positive surreal tolerance, choose $h$ smaller than that tolerance divided by a fixed real bound for $R_x$.

For the sign assertion write $x=r+\eps$ with $r\in[a,b]$. If $g(r)>0$, the standard part decides the sign. If $g(r)=0$ and the germ is not zero, factor it as $(X-r)^m u(X)$ with $u(r)\ne0$. At an interior zero, nonnegativity forces even $m$ and $u(r)>0$. At an endpoint, its leading term has the correct sign on the side allowed by $x\in[a,b]$. Substitution preserves these signs. A zero germ causes no difficulty.
\end{proof}
The uniqueness assertion is deliberately about the Taylor rule, not about arbitrary functions satisfying the same differential equation in the fine topology.

\section{Canonical sine, cosine, and finite-angle identities}\label{sec:finite}
\subsection{Definition without a divergent global series}
For $\theta=r+\eps\in\OO_{\R}$ define
\begin{align}
 \sin\theta&=\sin r\sum_{n\ge0}\frac{(-1)^n\eps^{2n}}{(2n)!}
 +\cos r\sum_{n\ge0}\frac{(-1)^n\eps^{2n+1}}{(2n+1)!},\\
 \cos\theta&=\cos r\sum_{n\ge0}\frac{(-1)^n\eps^{2n}}{(2n)!}
 -\sin r\sum_{n\ge0}\frac{(-1)^n\eps^{2n+1}}{(2n+1)!}.
 \label{eq:sincos}
\end{align}
These are exactly the scalar lifts of the ordinary functions. Set $\cis\theta=\cos\theta+i\sin\theta$. Unless an extension is explicitly named, lowercase sine, cosine, and tangent in this article have finite surreal real arguments.

\begin{theorem}[Elementary trigonometry on finite surreal angles]\label{thm:identities}
For $\alpha,\beta\in\OO_{\R}$,
\begin{align}
 \cos^2\alpha+\sin^2\alpha&=1,\\
 \sin(\alpha+\beta)&=\sin\alpha\cos\beta+\cos\alpha\sin\beta,\\
 \cos(\alpha+\beta)&=\cos\alpha\cos\beta-\sin\alpha\sin\beta,\\
 \cis(\alpha+\beta)&=\cis\alpha\,\cis\beta.
 \label{eq:addition}
\end{align}
Sine is odd, cosine is even, both have period $2\pi$, and
\[
 (\sin)'=\cos,\qquad (\cos)'=-\sin.
\]
For every ordinary integer $n$,
\begin{equation}\label{eq:demoivre}
 (\cis\alpha)^n=\cis(n\alpha).
\end{equation}
In particular, the usual double-angle, half-angle, sum-to-product, and product-to-sum identities hold, with every quotient restricted to a nonzero denominator.
\end{theorem}
\begin{proof}
Write $\alpha=r+\eps$ and $\beta=s+\eta$. The ordinary addition identities at $r,s$ and the formal identities for the power series in $\eps,\eta$ give \eqref{eq:addition}; all reorganizations are justified by \cref{lem:support}. The remaining identities follow by substitution, conjugation, differentiation, and ordinary finite induction. Negative integer powers use $\overline{\cis\alpha}=(\cis\alpha)^{-1}$.
\end{proof}

For convenient reference, several consequences are
\begin{align}
 \tan(\alpha+\beta)&=\frac{\tan\alpha+\tan\beta}{1-\tan\alpha\tan\beta},\\
 \sin\alpha+\sin\beta&=2\sin\frac{\alpha+\beta}{2}\cos\frac{\alpha-\beta}{2},\\
 \cos\alpha-\cos\beta&=-2\sin\frac{\alpha+\beta}{2}\sin\frac{\alpha-\beta}{2},\\
 2\sin\alpha\sin\beta&=\cos(\alpha-\beta)-\cos(\alpha+\beta),\\
 \cos(n\alpha)&=T_n(\cos\alpha),\\
 \sin(n\alpha)&=\sin\alpha\,U_{n-1}(\cos\alpha)\quad(n\ge1).
 \label{eq:chebyshev}
\end{align}
Here $T_0=1$, $T_1=X$, $T_{n+1}=2XT_n-T_{n-1}$ and $U_0=1$, $U_1=2X$, $U_{n+1}=2XU_n-U_{n-1}$. The last two identities follow from the same recurrence, so no new summability issue enters. The index $n$ is an ordinary finite integer, not an infinite surreal number.

\subsection{Order, monotonicity, and elementary inequalities}
\begin{proposition}[Signs and monotonicity]\label{prop:order}
Sine is positive on $(0,\pi)$ and negative on $(\pi,2\pi)$. Cosine is strictly decreasing on $[0,\pi]$, and sine is strictly increasing on $[-\pi/2,\pi/2]$. Moreover,
\begin{equation}\label{eq:jordan}
 \frac{2x}{\pi}\le\sin x\le x\qquad(0\le x\le\pi/2),
\end{equation}
and $|\sin x|\le|x|$ for every finite $x$.
\end{proposition}
\begin{proof}
An ordinary standard part strictly between $0$ and $\pi$ decides the positive sign. Near $0$, $\sin\eps=\eps(1+\text{infinitesimal})$; near $\pi$, $\sin(\pi-\eps)=\sin\eps$. These cover the remaining points of $(0,\pi)$. The other sign follows by translation by $\pi$.

For $0\le a<b\le\pi$, the identity
\[
 \cos a-\cos b=2\sin\frac{a+b}{2}\sin\frac{b-a}{2}>0
\]
proves strict decrease. The analogous sine difference uses a positive cosine of the midpoint on $[-\pi/2,\pi/2]$.

The ordinary real inequalities in \eqref{eq:jordan}, obtainable from concavity and the real derivative of sine, lift by \cref{prop:lift}. Finally, for $|x|\le1$ use the upper inequality and oddness; for $|x|\ge1$, use $|\sin x|\le1$ from the Pythagorean identity.
\end{proof}
This proof does not infer global monotonicity from the sign of a fine derivative. Such an inference would require additional hypotheses on a general surreal function.

\begin{proposition}[Exact infinitesimal leading orders]\label{prop:leading}
For $0\ne h\in\mm_{\R}$,
\begin{align}
 \sin h&=h-\frac{h^3}{6}+\frac{h^5}{120}+O(h^7),\\
 1-\cos h&=\frac{h^2}{2}-\frac{h^4}{24}+O(h^6),\\
 \tan h&=h+\frac{h^3}{3}+\frac{2h^5}{15}+O(h^7),
\end{align}
and, exactly,
\begin{equation}\label{eq:sinval}
 v(\sin h)=v(h),\qquad
 v(1-\cos h)=2v(h),\qquad v(\tan h)=v(h).
\end{equation}
\end{proposition}
\begin{proof}
The displayed truncations come from the strongly summable formal series. Dividing the three expressions by $h,h^2,h$, respectively, gives finite elements with nonzero standard parts $1,1/2,1$. Such units have valuation zero.
\end{proof}

\section{The full unit circle and the group of angles}\label{sec:circle}
\subsection{Infinitesimal logarithms}
For $h,\eta\in\mm_{\C}$ put
\begin{equation}\label{eq:explog}
 \Ex(h)=\sum_{n\ge0}\frac{h^n}{n!},\qquad
 \Logm(1+\eta)=\sum_{n\ge1}\frac{(-1)^{n+1}\eta^n}{n}.
\end{equation}
Formal power-series identities and \cref{lem:support} prove that these maps are mutually inverse between $\mm_{\C}$ and $1+\mm_{\C}$, that $\Ex(h+k)=\Ex(h)\Ex(k)$, and that the logarithm turns products into sums. Both commute with conjugation. In particular,
\begin{equation}\label{eq:logleading}
 v\bigl(\Logm(1+\eta)\bigr)=v(\eta)
 \quad(\eta\ne0),\qquad \cis h=\Ex(ih)\quad(h\in\mm_{\R}).
\end{equation}

Define the unit circle
\[
 \TT(\No)=\{u\in\SC:u\bar u=1\}.
\]
Its elements all have finite coordinates, even though the circle is a proper class. Let $\TT(\R)$ denote the ordinary complex unit circle.

\begin{theorem}[Polar decomposition and the angle group]\label{thm:polar}
The map
\[
 \cis:(\OO_{\R},+)\longrightarrow(\TT(\No),\cdot)
\]
is surjective and has kernel $2\pi\Z$. Thus
\begin{equation}\label{eq:anglegroup}
 \OO_{\R}/2\pi\Z\ \cong\ \TT(\No).
\end{equation}
More precisely, there is a canonical group isomorphism
\begin{equation}\label{eq:circlesplit}
 \TT(\No)\longrightarrow\TT(\R)\times(\mm_{\R},+),\qquad
 u\longmapsto\left(u_0,-i\Logm(u/u_0)\right),\quad u_0=\st(u).
\end{equation}
Every $z\ne0$ therefore has a unique modulus $r=|z|>0$ and a unique angle class $[\theta]\in\OO_{\R}/2\pi\Z$ such that $z=r\cis\theta$.
\end{theorem}
\begin{proof}
Taking standard parts in $u\bar u=1$ gives $u_0\in\TT(\R)$. Put $q=u/u_0\in1+\mm_{\C}$. Since $\bar q=q^{-1}$,
\[
 \overline{\Logm q}=\Logm\bar q=\Logm(q^{-1})=-\Logm q.
\]
Thus $\Logm q=i\delta$ for a unique $\delta\in\mm_{\R}$. Choose any ordinary argument $\theta_0$ of $u_0$. Then
\[
 u=u_0\Ex(i\delta)=\cis(\theta_0+\delta).
\]
This proves surjectivity. The standard-part map and the logarithmic group law prove that \eqref{eq:circlesplit} is a homomorphism with inverse $(u_0,\delta)\mapsto u_0\Ex(i\delta)$.

If $\cis(r+\eps)=1$, standard part gives $r\in2\pi\Z$. It follows that $\Ex(i\eps)=1$, and logarithmic inversion gives $\eps=0$. This proves the kernel statement. Apply the circle result to $u=z/|z|$ for polar decomposition.
\end{proof}

The infinitesimal component in \eqref{eq:circlesplit} is independent of an ordinary argument branch. A numerical representative of an angle class requires a cut. Each class has exactly one representative in $[0,2\pi)$, obtained by adding or subtracting an ordinary multiple of $2\pi$ and checking the actual surreal inequalities. Near the endpoints, the infinitesimal correction matters. For example, the representative of $-\eps$ with $\eps>0$ infinitesimal is $2\pi-\eps$, not an ordinary angle.

\begin{remark}[Two different branch conventions]
A representative whose \emph{actual value} lies in $(-\pi,\pi]$ has a branch jump at the negative real axis, including across infinitesimal perturbations. One may instead fix an argument of the \emph{standard direction} and retain its infinitesimal correction without reducing it again. The latter convention is fine-locally analytic but can lie infinitesimally outside $(-\pi,\pi]$. The global logarithm in the supplied analysis manuscript uses this second kind of convention. Angle classes avoid the ambiguity entirely.
\end{remark}

\subsection{Rational coordinates and angle division}
\begin{theorem}[Half-angle parametrization of every surreal direction]\label{thm:cayley}
The map
\begin{equation}\label{eq:cayley}
 u(t)=\frac{1+it}{1-it}
 =\frac{1-t^2}{1+t^2}+i\frac{2t}{1+t^2},\qquad t\in\No,
\end{equation}
is a bijection from $\No$ to $\TT(\No)\setminus\{-1\}$. Its inverse is
\[
 t=\frac{\im u}{1+\re u}.
\]
If $u=\cis\theta\ne-1$, this coordinate is $\tan(\theta/2)$ for the representative $-\pi<\theta<\pi$. Infinite $t$ are allowed and describe directions infinitesimally close to $-1$.
\end{theorem}
\begin{proof}
The displayed rational expression has norm one and cannot equal $-1$, because $1+\re u(t)=2/(1+t^2)>0$. Conversely, if $u=c+is$ has $c^2+s^2=1$ and $c\ne-1$, then $1+c>0$. Substituting $t=s/(1+c)$ in \eqref{eq:cayley} recovers $c+is$. The half-angle identity follows from \cref{thm:identities} and the chosen interval.
\end{proof}

\begin{corollary}[Roots and regular polygons]\label{cor:rootsunity}
For an ordinary $n\ge1$, the equation $v^n=\cis\theta$ on $\TT(\No)$ has exactly $n$ solutions,
\[
 v_k=\cis\frac{\theta+2\pi k}{n},\qquad 0\le k<n.
\]
In particular, the $n$th roots of unity are exactly the ordinary complex roots of unity. For $n\ge3$, a regular $n$-gon of radius $r\in\No_{>0}$ has side length $2r\sin(\pi/n)$.
\end{corollary}
\begin{proof}
The formulas give distinct solutions by the kernel in \cref{thm:polar}. Any other solution has some finite angle $\phi$, and $n\phi-\theta\in2\pi\Z$, which reduces to one of the listed classes. The chord formula follows either from \eqref{eq:gram} or from the cosine double-angle identity.
\end{proof}
The uniqueness of an infinitesimal $n$th root of a direction near $1$ is also visible in \eqref{eq:circlesplit}: its infinitesimal angle is divided by $n$.

\section{Inverse functions and the angular metric}\label{sec:inverse}
\subsection{Inverse trigonometry on its full natural surreal domains}
\begin{theorem}[Inverse tangent, sine, and cosine]\label{thm:inverse}
There are unique inverse functions
\begin{align*}
 \arctan &: \No\longrightarrow(-\pi/2,\pi/2),\\
 \arcsin &: [-1,1]\longrightarrow[-\pi/2,\pi/2],\\
 \arccos &: [-1,1]\longrightarrow[0,\pi],
\end{align*}
with the usual inverse identities. All values are finite surreal angles. They satisfy
\begin{align}
 \cos(\arctan t)&=\frac1{\sqrt{1+t^2}},&
 \sin(\arctan t)&=\frac{t}{\sqrt{1+t^2}},\\
 \arccos t&=\frac\pi2-\arcsin t.
\end{align}
At every point of their open domains, including infinite $t$ for inverse tangent and points infinitesimally close to $\pm1$ for the other two,
\begin{equation}\label{eq:inversederivatives}
 (\arctan)'(t)=\frac1{1+t^2},\quad
 (\arcsin)'(t)=\frac1{\sqrt{1-t^2}},\quad
 (\arccos)'(t)=-\frac1{\sqrt{1-t^2}}.
\end{equation}
\end{theorem}
\begin{proof}
The unit direction $(1+it)/\sqrt{1+t^2}$ has positive real part, so \cref{thm:polar} and the sign rules give a unique angle in $(-\pi/2,\pi/2)$. This defines inverse tangent and proves its two coordinate formulas. Similarly, the directions
\[
 \sqrt{1-t^2}+it,\qquad t+i\sqrt{1-t^2}
\]
give inverse sine and inverse cosine on their stated ranges. Strict monotonicity in \cref{prop:order} proves uniqueness, and direct substitution proves the inverse identities.

Locally, square roots of nonzero positive elements have their binomial expansions after division by the base value. A local argument difference is the imaginary part of the infinitesimal logarithm of a direction ratio. These two facts give fine-local analytic expansions of the three inverse functions at every point under discussion. Differentiating the inverse identities then gives \eqref{eq:inversederivatives}; the relevant cosine or sine is nonzero. At an infinite inverse-tangent input, the same computation is valid after taking increments small relative to $1+it$. No finite bound on $t$ is used.
\end{proof}

For positive infinite $t$, and hence infinitesimal $1/t$,
\begin{equation}\label{eq:ataninfty}
 \arctan t=\frac\pi2-\arctan(1/t)
 =\frac\pi2-\frac1t+\frac1{3t^3}-\frac1{5t^5}+O(t^{-7}).
\end{equation}
For negative $t$ use oddness. Thus an infinite slope corresponds to a finite angle infinitesimally short of a right angle.

\begin{proposition}[Endpoint ramification of inverse cosine]\label{prop:acosendpoint}
For $0<\delta\in\mm_{\R}$,
\begin{align}
 \arccos(1-\delta)
 &=2\arcsin\sqrt{\delta/2}\label{eq:acosseries}\\
 &=\sqrt{2\delta}\sum_{n\ge0}
 \frac{\binom{2n}{n}}{8^n(2n+1)}\delta^n\notag\\
 &=\sqrt{2\delta}\left(1+\frac\delta{12}
       +\frac{3\delta^2}{160}+\frac{5\delta^3}{896}+O(\delta^4)\right).
 \notag
\end{align}
The series is strongly summable after adjoining $\sqrt\delta$, and
\begin{equation}\label{eq:acosval}
 v\bigl(\arccos(1-\delta)\bigr)=\tfrac12v(\delta).
\end{equation}
\end{proposition}
\begin{proof}
The double-angle identity gives $\cos(2\arcsin\sqrt{\delta/2})=1-\delta$. The angle is positive infinitesimal, so inverse uniqueness gives the first equality. Substitute the ordinary formal series for inverse sine. Every remaining factor in parentheses has standard part one, proving the valuation formula.
\end{proof}
The functions inverse sine and inverse tangent are valuation-preserving at zero, whereas inverse cosine at $1$ is a square-root branch. This distinction will govern tangency.

\subsection{Chord distance and infinitesimal angular accuracy}
For $u,v\in\TT(\No)$, let $d_\theta(u,v)\in[0,\pi]$ be the least absolute representative of their angle difference. Equivalently,
\[
 d_\theta(u,v)=\arccos\re(\bar uv).
\]
This is an angular metric: choose shortest representatives for two successive differences, add them, and reduce modulo $2\pi$; reduction cannot increase the least absolute value, which proves the triangle inequality.

\begin{theorem}[Angular and chordal comparison]\label{thm:metric}
For $u,v\in\TT(\No)$, with $d=d_\theta(u,v)$,
\begin{equation}\label{eq:metricbounds}
 |u-v|=2\sin(d/2),\qquad
 |u-v|\le d\le\frac\pi2|u-v|.
\end{equation}
For $u\ne v$,
\begin{equation}\label{eq:metricval}
 v(d_\theta(u,v))=v(u-v).
\end{equation}
If $\st(u)=\st(v)$, their unique infinitesimal oriented angle difference is
\begin{equation}\label{eq:localangle}
 \delta=-i\Logm(v/u),\qquad
 v(\delta)=v(v-u),\qquad d_\theta(u,v)=|\delta|.
\end{equation}
\end{theorem}
\begin{proof}
Expanding $|u-v|^2=2-2\re(\bar uv)$ and using the half-angle identity proves the first equality in \eqref{eq:metricbounds}. Jordan's inequality \eqref{eq:jordan} proves both bounds. Multiplication by a nonzero ordinary real constant does not change valuation, so these bounds force \eqref{eq:metricval}. Finally $v/u\in1+\mm_{\C}$, and \eqref{eq:logleading} and the circle logarithm prove \eqref{eq:localangle}.
\end{proof}
Thus angular precision is exactly the precision of normalized directions, not merely comparable up to an unspecified non-Archimedean error.

\begin{corollary}[Stability of direction under a relative perturbation]\label{cor:directionstability}
If $z\ne0$ and $h/z=\eta\in\mm_{\C}$, then a local angle difference between $z+h$ and $z$ is
\[
 \Delta\theta=\im\Logm(1+\eta)
 =\im\eta-\tfrac12\im(\eta^2)+O(|\eta|^3).
\]
In particular, $v(\Delta\theta)\ge v(\eta)$ when the difference is nonzero. Equality holds if the leading coefficient of $\eta$ has nonzero imaginary part.
\end{corollary}
\begin{proof}
The real part of the logarithm accounts for the change in modulus and its imaginary part for the change in direction. Taking imaginary parts cannot lower a Hahn valuation. If the first coefficient has nonzero imaginary part, the first term cannot cancel with the higher powers.
\end{proof}

\section{Triangles at arbitrary surreal scales}\label{sec:triangles}
\subsection{Angles and the exact angle sum}
Let $A,B,C\in\SC$ be noncollinear, with side lengths
\[
 a=|B-C|,\qquad b=|C-A|,\qquad c=|A-B|,
\]
and positive area
\[
 \Delta=\tfrac12|[B-A,C-A]|.
\]
All four quantities are positive surreals; no finiteness assumption is made. Define the interior angle $\alpha$ at $A$ by
\begin{equation}\label{eq:interiorangle}
 \cos\alpha=\frac{\langle B-A,C-A\rangle}{bc},\qquad
 \sin\alpha=\frac{2\Delta}{bc},\qquad 0<\alpha<\pi.
\end{equation}
Equation~\eqref{eq:gram} shows that the right sides lie on the upper unit semicircle, so \cref{thm:polar} gives existence and uniqueness. Define $\beta,\gamma$ cyclically.

\begin{theorem}[Euclidean angle sum over the surreal field]\label{thm:anglesum}
Every noncollinear surcomplex triangle satisfies
\[
 \alpha+\beta+\gamma=\pi.
\]
\end{theorem}
\begin{proof}
A translation, rotation, and possibly conjugation put $A=0$, $B=c>0$, $C=x+iy$ with $y>0$. Let $\alpha=\Arg C\in(0,\pi)$ and $\delta=\Arg(C-c)\in(0,\pi)$. Since
\[
 \im\bigl(\bar C(C-c)\bigr)=cy>0,
\]
we have $\sin(\delta-\alpha)>0$. The difference lies in $(-\pi,\pi)$, so $\delta>\alpha$. The angle at $B$ is $\pi-\delta$, and the angle at $C$ is $\delta-\alpha$; these are both in $(0,\pi)$. Their sum with $\alpha$ is $\pi$.
\end{proof}
The proof uses order, angle classes, and determinants. It applies equally to an infinitesimal triangle, an infinite triangle, or a triangle with sides on different scales.

\subsection{Sine, cosine, and tangent laws}
\begin{theorem}[Triangle laws]\label{thm:trianglelaws}
For every noncollinear surcomplex triangle,
\begin{align}
 a^2&=b^2+c^2-2bc\cos\alpha,\\
 \frac a{\sin\alpha}&=\frac b{\sin\beta}=\frac c{\sin\gamma}
 =\frac{abc}{2\Delta}=2R,\\
 \Delta&=\tfrac12bc\sin\alpha,\\
 \frac{a-b}{a+b}
 &=\frac{\tan((\alpha-\beta)/2)}{\tan((\alpha+\beta)/2)}.
 \label{eq:trianglelaws}
\end{align}
Here $R$ is the radius of the unique circumcircle. All denominators in these formulas are nonzero.
\end{theorem}
\begin{proof}
Expand $|(C-A)-(B-A)|^2$ to obtain the cosine law. Equation~\eqref{eq:interiorangle} and its cyclic versions give the area identity and the common sine-law value.

For the circumcircle, again set $A=0$, $B=c$, $C=x+iy$, $y>0$. The equations for a center $O=X+iY$ equidistant from the three vertices are
\[
 2cX=c^2,\qquad 2xX+2yY=x^2+y^2=b^2.
\]
They have the unique solution $X=c/2$, $Y=(b^2-cx)/(2y)$. A direct expansion, using $a^2=(x-c)^2+y^2$, gives
\[
 |O|^2=\frac{a^2b^2}{4y^2}.
\]
Hence $R=ab/(2y)=abc/(4\Delta)$, as claimed. This also proves existence without a compactness argument.

Finally the sine law changes $(a-b)/(a+b)$ into
$(\sin\alpha-\sin\beta)/(\sin\alpha+\sin\beta)$. Apply the sum-to-product identities. Since $0<\alpha+\beta<\pi$ and $|\alpha-\beta|<\pi$, the cancellations and tangents are legitimate.
\end{proof}

\begin{corollary}[Right triangles and the full range of slopes]\label{cor:right}
A triangle is right at $C$ exactly when $c^2=a^2+b^2$. In a right triangle with positive legs $p,q$, its acute angle opposite $q$ is $\arctan(q/p)$. The ratio $q/p$ can be any positive surreal, including an infinite one.
\end{corollary}
\begin{proof}
The cosine law and the unique zero of cosine on $(0,\pi)$ prove the first assertion. The second follows from the defining coordinate formulas of inverse tangent.
\end{proof}

\begin{theorem}[Existence from side lengths and similarity]\label{thm:sss}
Positive surreals $a,b,c$ are the side lengths of a noncollinear surcomplex triangle if and only if all three strict triangle inequalities hold. The triangle is unique up to a Euclidean isometry and reflection. Triangles with the same three angles have proportional corresponding sides.
\end{theorem}
\begin{proof}
Necessity is the norm triangle inequality, with equality excluded by noncollinearity. For sufficiency place $A=0$, $B=c$, and set
\[
 x=\frac{b^2+c^2-a^2}{2c},\qquad y=\sqrt{b^2-x^2}.
\]
The identity
\begin{equation}\label{eq:heronfactor}
 4c^2(b^2-x^2)
 =(a+b+c)(-a+b+c)(a-b+c)(a+b-c)
\end{equation}
shows that $y>0$. Then $C=x+iy$ has the required distances. The same equations force $x$ and $|y|$, proving uniqueness. The sine law gives proportional sides for equal angles.
\end{proof}

\subsection{Area, half-angles, and the two radii}
Put $s=(a+b+c)/2$ and let $r$ denote the inradius.

\begin{theorem}[Heron and half-angle formulas]\label{thm:heron}
One has
\begin{align}
 \Delta^2&=s(s-a)(s-b)(s-c),& r&=\frac\Delta s,\\
 \sin\frac\alpha2&=\sqrt{\frac{(s-b)(s-c)}{bc}},&
 \cos\frac\alpha2&=\sqrt{\frac{s(s-a)}{bc}},\\
 \tan\frac\alpha2&=\frac r{s-a}.
 \label{eq:halfangletriangle}
\end{align}
The square roots are the positive ones. The incenter is
\begin{equation}\label{eq:incenter}
 I=\frac{aA+bB+cC}{2s}.
\end{equation}
\end{theorem}
\begin{proof}
With the coordinates in \cref{thm:sss}, $\Delta=cy/2$, so \eqref{eq:heronfactor} proves Heron's formula. The cosine law and $\sin^2(\alpha/2)=(1-\cos\alpha)/2$ give the first half-angle formula; the other follows by addition to one. Positivity fixes both square roots.

The point in \eqref{eq:incenter} lies strictly inside the triangle because its barycentric weights are positive. Its distance to $AB$ is $c/(2s)$ times the altitude from $C$, hence $\Delta/s$; the other two distances agree by cyclic symmetry. Thus it is the incenter and $r=\Delta/s$. Dividing the half-angle formulas and using Heron's identity gives the tangent formula.
\end{proof}

\begin{theorem}[Euler's incenter--circumcenter identity]\label{thm:euler}
If $O$ is the circumcenter and $I$ the incenter, then
\begin{equation}\label{eq:euler}
 |O-I|^2=R(R-2r).
\end{equation}
Consequently $R\ge2r$, with equality exactly for an equilateral triangle.
\end{theorem}
\begin{proof}
Translate $O$ to zero, so $|A|=|B|=|C|=R$. For arbitrary positive weights $w_j$ and vectors $Z_j$, expansion of the inner products gives
\[
 \left|\frac{\sum_jw_jZ_j}{\sum_jw_j}\right|^2
 =\frac{\sum_jw_j|Z_j|^2}{\sum_jw_j}
 -\frac{\sum_{j<k}w_jw_k|Z_j-Z_k|^2}{(\sum_jw_j)^2}.
\]
Apply this with $(w_1,w_2,w_3)=(a,b,c)$. The numerator in the second term is
\[
 ab\,c^2+ac\,b^2+bc\,a^2=2sabc.
\]
Hence $|I|^2=R^2-abc/(2s)=R^2-2Rr$. Positivity proves the inequality. If equality holds, $I=O$. Every side then has the same positive distance $r$ from the center of the circle, so its length is $2\sqrt{R^2-r^2}$; the triangle is equilateral. The converse is immediate.
\end{proof}
No limiting or Archimedean argument is hidden in any of these identities.

\section{Bisectors, concurrency, and cyclic quadrilaterals}\label{sec:incidence}
\subsection{The angle-bisector theorem and trigonometric Ceva}
\begin{theorem}[Sine ratios for a cevian]\label{thm:cevian}
Let $D$ lie strictly inside the side $BC$ of a noncollinear triangle. Then
\begin{equation}\label{eq:cevian}
 \frac{|B-D|}{|D-C|}
 =\frac{c\sin\angle BAD}{b\sin\angle DAC}.
\end{equation}
In particular, the internal bisector of $\alpha$ meets $BC$ at the unique point satisfying $|B-D|/|D-C|=c/b$.
\end{theorem}
\begin{proof}
The two triangles $ABD$ and $ADC$ share an altitude to the line $BC$, so their areas are in the ratio $|B-D|/|D-C|$. Compute the same areas using the common side $AD$ and their included angles at $A$. This gives \eqref{eq:cevian}. For the bisector, equal sines of equal positive half-angles cancel. Conversely, a ray of angle $\alpha/2$ lies inside the angle at $A$ and meets $BC$; or one may construct the point of ratio $c/b$ and use the sine ratio together with the strict increase of $\sin x/\sin(\alpha-x)$ on $(0,\alpha)$. The latter increase follows directly from the sine subtraction identity, whose numerator is $\sin\alpha\sin(y-x)>0$ for $x<y$.
\end{proof}

\begin{theorem}[Trigonometric Ceva]\label{thm:ceva}
Let $D,E,F$ lie strictly inside $BC,CA,AB$, respectively. The lines $AD,BE,CF$ are concurrent if and only if
\begin{equation}\label{eq:trigceva}
 \frac{\sin\angle BAD}{\sin\angle DAC}
 \frac{\sin\angle CBE}{\sin\angle EBA}
 \frac{\sin\angle ACF}{\sin\angle FCB}=1.
\end{equation}
\end{theorem}
\begin{proof}
For an interior point with positive barycentric coordinates $(p_A,p_B,p_C)$, its trace on $BC$ has ratio $|B-D|/|D-C|=p_C/p_B$. The two cyclic ratios are $p_A/p_C$ and $p_B/p_A$, so their product is one. Conversely, three prescribed positive ratios with product one determine such positive coordinates, hence a common point. This proves the ordinary ratio form of Ceva by field algebra alone. Substitute \eqref{eq:cevian} and its cyclic versions. The factors $c/b$, $a/c$, and $b/a$ cancel, giving \eqref{eq:trigceva}.
\end{proof}
The restriction to interior traces avoids directed-length conventions. A directed Menelaus or Ceva formula can be obtained in the same way after fixing signs, but is not needed for the present development.

\subsection{Chords, inscribed angles, and Ptolemy}
If $z_j=O+R\cis\theta_j$ lie on a circle, then
\begin{equation}\label{eq:chordfactor}
 z_2-z_1=2iR\cis\frac{\theta_1+\theta_2}{2}
                 \sin\frac{\theta_2-\theta_1}{2}.
\end{equation}
The formula follows by subtracting two unit phases and applying the addition identities. It holds at every surreal radius.

\begin{proposition}[Inscribed-angle theorem]\label{prop:inscribed}
For three distinct points $z_1,z_2,z_3$ on a circle, the directed angle from $z_1-z_3$ to $z_2-z_3$, taken modulo $\pi$, is the half of the directed central angle from $z_1-O$ to $z_2-O$. Here halving sends a central angle modulo $2\pi$ to an angle modulo $\pi$. In particular, points on a common arc see a fixed chord under equal interior angles, and an angle subtending a diameter is right.
\end{proposition}
\begin{proof}
Apply \eqref{eq:chordfactor} to both vectors based at $z_3$. Their quotient has phase $(\theta_2-\theta_1)/2$ times a nonzero real scalar. A real sign affects an argument only by $\pi$, proving the directed statement. Choosing the actual interior-angle ranges gives the stated consequences.
\end{proof}

\begin{theorem}[Ptolemy inequality and its cyclic equality]\label{thm:ptolemy}
For arbitrary four surcomplex points,
\begin{equation}\label{eq:ptolemy}
 |z_1-z_3|\,|z_2-z_4|
 \le |z_1-z_2|\,|z_3-z_4|+|z_1-z_4|\,|z_2-z_3|.
\end{equation}
If the four points are distinct and cyclically ordered on a circle, equality holds.
\end{theorem}
\begin{proof}
The algebraic identity
\[
 (z_1-z_3)(z_2-z_4)
 =(z_1-z_2)(z_3-z_4)+(z_1-z_4)(z_2-z_3)
\]
and the norm triangle inequality prove \eqref{eq:ptolemy}.

For cyclic points choose finite lifts
\[
 \theta_1<\theta_2<\theta_3<\theta_4<\theta_1+2\pi.
\]
By \eqref{eq:chordfactor}, the two products on the right of the algebraic identity have the same unit phase and strictly positive real multiples of that phase: all the half-difference sines involved have the prescribed sign. Equality therefore holds in the norm triangle inequality.
\end{proof}
This illustrates a general principle: algebraic identities among finitely many vectors often survive unchanged, while any analytic interpretation of the associated angles must still be supplied.

\section{Valuation geometry and extremely thin triangles}\label{sec:valuationgeometry}
\subsection{An exact dictionary between sides, angles, and area}
For an interior angle $\alpha\in(0,\pi)$ define its endpoint distance
\[
 d(\alpha)=\min\{\alpha,\pi-\alpha\}\in(0,\pi/2].
\]
It measures proximity to either a zero angle or a straight angle.

\begin{theorem}[Valuation form of the triangle laws]\label{thm:valuationtriangle}
For a noncollinear surcomplex triangle,
\begin{align}
 v(\sin\alpha)&=v(d(\alpha)),\\
 v(a)-v(b)&=v(d(\alpha))-v(d(\beta)),\\
 v(\Delta)&=v(b)+v(c)+v(d(\alpha)),\\
 2v(\Delta)&=v(s)+v(s-a)+v(s-b)+v(s-c),\\
 v(R)&=v(a)+v(b)+v(c)-v(\Delta),\\
 v(r)&=v(\Delta)-v(s).
 \label{eq:valuationtriangle}
\end{align}
The minimum of $v(a),v(b),v(c)$ is attained at least twice.
\end{theorem}
\begin{proof}
Jordan's inequality compares $\sin d(\alpha)$ to $d(\alpha)$ by nonzero ordinary constants, so their valuations agree. Since $\sin\alpha=\sin d(\alpha)$, the first assertion follows. Apply valuation to the sine law, area formula, Heron's identity, and radius formulas to obtain the others.

If, for example, $v(a)<v(b),v(c)$, then $b/a$ and $c/a$ are positive infinitesimals. Their sum is less than one, contradicting $a<b+c$. Thus no side can be the unique dominant valuation scale.
\end{proof}
The last statement is a non-Archimedean form of the triangle inequality. It does not say that all three sides have equal valuation: a very short third side is possible.

Choose a longest side and scale it to length one. The remaining vertex can be written as $x+iy$, $y>0$, with $0<x<1$ and both coordinates finite. The normalized real shadow is noncollinear exactly when $y$ is appreciable. If $y$ is infinitesimal, the real shadow loses information that remains visible in \eqref{eq:valuationtriangle}. Vertices may also coalesce in the shadow; the next theorem treats the clean flat case where they do not.

\subsection{A flat-triangle theorem with exact leading constants}
\begin{theorem}[Quadratic slack and reciprocal circumradius]\label{thm:flat}
Let
\[
 A=0,\qquad B=1,\qquad C=x+iy,
\]
where $x$ is finite with $0<\st(x)<1$, and $0<y\in\mm_{\R}$. Put
\[
 a=\sqrt{(1-x)^2+y^2},\qquad b=\sqrt{x^2+y^2},\qquad c=1.
\]
Then the two small angles and the large angle are exactly
\begin{equation}\label{eq:flatangles}
 \alpha=\arctan(y/x),\quad
 \beta=\arctan(y/(1-x)),\quad
 \gamma=\pi-\alpha-\beta.
\end{equation}
Moreover,
\begin{align}
 \alpha&=\frac yx-\frac{y^3}{3x^3}+O(y^5),&
 \beta&=\frac y{1-x}-\frac{y^3}{3(1-x)^3}+O(y^5),\\
 a+b-1&=\frac{y^2}{2x(1-x)}
 -\frac{y^4}{8}\left(\frac1{x^3}+\frac1{(1-x)^3}\right)+O(y^6),
 \label{eq:flatslack}\\
 R&=\frac{x(1-x)}{2y}
 \left[1+\frac{y^2}{2}\left(\frac1{x^2}+\frac1{(1-x)^2}\right)+O(y^4)\right],\\
 r&=\frac y2\left[1-\frac{y^2}{4x(1-x)}+O(y^4)\right].
\end{align}
Consequently,
\begin{align}
 v(\alpha)=v(\beta)=v(\pi-\gamma)&=v(y),\\
 v(a+b-1)&=2v(y),& v(R)&=-v(y),&v(r)&=v(y).
 \label{eq:flatvals}
\end{align}
\end{theorem}
\begin{proof}
Both $x$ and $1-x$ are positive appreciable, so the base angles are acute and their slopes give \eqref{eq:flatangles}. The inverse-tangent series gives their expansions. Expand
\[
 b=x\sqrt{1+(y/x)^2},\qquad
 a=(1-x)\sqrt{1+(y/(1-x))^2}
\]
by the binomial series, which is legitimate at these infinitesimal arguments. This gives \eqref{eq:flatslack}. Since $\Delta=y/2$, the exact radius formulas are
\[
 R=\frac{ab}{2y},\qquad r=\frac{y}{a+b+1}.
\]
Multiplication and geometric inversion yield the displayed expansions. Each indicated leading coefficient is positive appreciable, including the sum of the two leading angle coefficients. Taking valuations gives \eqref{eq:flatvals}.
\end{proof}
The theorem quantifies a phenomenon invisible in the ordinary real shadow: the two small angles are first-order in height, the excess in the triangle inequality is second-order, and the circumradius is inverse-order.

\begin{example}[A symmetric infinitesimal-height triangle]\label{ex:flat}
Take $x=1/2$ and $y=\eps>0$ infinitesimal. Then
\begin{align*}
 a=b&=\tfrac12\sqrt{1+4\eps^2},\\
 \alpha=\beta&=\arctan(2\eps),\\
 R&=\frac1{8\eps}+\frac\eps2,\\
 r&=\frac{\eps}{1+\sqrt{1+4\eps^2}}
   =\frac\eps2-\frac{\eps^3}{2}+\eps^5+O(\eps^7).
\end{align*}
The circumradius is infinite although all vertices lie within a bounded ordinary region. Its existence is not contradicted by the fact that the standard-part points are collinear.
\end{example}

\begin{example}[An infinite triangle of ordinary finite area]\label{ex:infinite}
Let $L>0$ be infinite and take $A=0$, $B=L$, $C=L+i/L$. The legs at $B$ are $L$ and $L^{-1}$, while
\[
 \Delta=\frac12,\qquad
 \alpha=\arctan(L^{-2}),\quad\beta=\frac\pi2,\quad
 \gamma=\frac\pi2-\alpha.
\]
Its circumradius is $\tfrac12\sqrt{L^2+L^{-2}}$, and its inradius is
\[
 r=\frac{L+L^{-1}-\sqrt{L^2+L^{-2}}}{2}
  =\frac1{2L}-\frac1{4L^3}+O(L^{-7}).
\]
Thus an infinite side, an infinitesimal side, an ordinary area, and an infinitesimal angle coexist without changing any triangle identity.
\end{example}

\section{Trigonometric intersections and infinitesimal contact}\label{sec:contact}
\subsection{A complete first-harmonic intersection theorem}
\begin{theorem}[Amplitude--phase reduction and intersection count]\label{thm:amplitude}
Let $A,B,D\in\No$ with $(A,B)\ne(0,0)$, and put $\rho=\sqrt{A^2+B^2}>0$. Choose the finite angle class $\phi$ satisfying
\[
 \cos\phi=A/\rho,\qquad \sin\phi=B/\rho.
\]
Then
\begin{equation}\label{eq:amplitude}
 A\cos\theta+B\sin\theta=\rho\cos(\theta-\phi).
\end{equation}
The equation $A\cos\theta+B\sin\theta=D$ has, modulo $2\pi$, no solutions if $|D|>\rho$, one solution if $|D|=\rho$, and exactly two solutions if $|D|<\rho$. In the last case they are
\[
 \theta=\phi\pm\arccos(D/\rho).
\]
\end{theorem}
\begin{proof}
Polar decomposition gives $\phi$, and the cosine subtraction identity proves \eqref{eq:amplitude}. The range, strict monotonicity, and inverse cosine on $[0,\pi]$ give the count, including the two endpoints. The two signs coincide modulo $2\pi$ exactly at those endpoints.
\end{proof}
Geometrically, this is the complete intersection classification of a line $AX+BY=D$ with the unit circle. Neither an infinite normal vector nor an infinitesimal distance from tangency creates an exception.

\subsection{Square-root splitting and its conditioning}
\begin{theorem}[Tangency, angular splitting, and loss of valuation]\label{thm:tangency}
In \cref{thm:amplitude}, suppose $D=\rho(1-\delta)$ with $0<\delta\in\mm_{\R}$. The two solutions near $\phi$ are
\begin{equation}\label{eq:tangency}
 \theta_\pm=\phi\pm\sqrt{2\delta}
 \left(1+\frac\delta{12}+\frac{3\delta^2}{160}+O(\delta^3)\right).
\end{equation}
Their separation has valuation $\tfrac12v(\delta)$. The derivatives of the left side of the equation at these roots have absolute value
\begin{equation}\label{eq:jacobiancontact}
 \rho\sqrt{\delta(2-\delta)}.
\end{equation}
If $e\ne0$ and $e\prec\delta$, set $q(\delta)=\arccos(1-\delta)$. Then $\delta+e>0$ and
\begin{equation}\label{eq:contactstability}
 q(\delta+e)-q(\delta)
 =\frac{e}{\sqrt{\delta(2-\delta)}}
       \bigl(1+O(e/\delta)\bigr),
\end{equation}
where the correction factor has standard part one. In particular,
\begin{equation}\label{eq:contactval}
 v\bigl(q(\delta+e)-q(\delta)\bigr)
 =v(e)-\tfrac12v(\delta).
\end{equation}
\end{theorem}
\begin{proof}
Equation~\eqref{eq:tangency} is \cref{prop:acosendpoint}. The separation is $2q(\delta)$, so its valuation is unchanged by the ordinary factor two. Differentiation of \eqref{eq:amplitude} gives magnitude $\rho\sin q(\delta)$, and
\[
 \sin q(\delta)=\sqrt{1-(1-\delta)^2}=\sqrt{\delta(2-\delta)}.
\]
For the perturbation formula, write $e=\delta u$ with $u$ infinitesimal. The exact expression $q(\delta)=\sqrt{2\delta}\,H(\delta)$ in \eqref{eq:acosseries}, where $H(0)=1$, may be expanded in the two infinitesimals $\delta,u$. Its first derivative in $e$ is $1/\sqrt{\delta(2-\delta)}$. Every later term, relative to that first term, contains a factor $u=e/\delta$ times a finite factor. This proves \eqref{eq:contactstability} by support-controlled substitution. Its first term is nonzero and the parenthesized multiplier is a unit of standard part one, proving \eqref{eq:contactval}.
\end{proof}
The loss of one half of $v(\delta)$ is the inverse-Jacobian loss for a nearly double root. The condition $e\prec\delta$ cannot be omitted from a uniform branch-matching theorem: $e=-\delta$ merges the roots, and $e<-\delta$ destroys the two real intersections.

\begin{example}[An infinite tangent coordinate]
The line $X=-1+\delta$ meets the circle near $-1$. Its two points have $Y=\pm\sqrt{\delta(2-\delta)}$. Their half-angle coordinates from \eqref{eq:cayley} are
\[
 t=\pm\sqrt{\frac{2-\delta}{\delta}},
\]
which are infinite, even though the points and their angles are finite. Rational circle coordinates and finite angles therefore complement each other rather than having identical size behavior.
\end{example}

\section{Canonical complex trigonometry on a class-sized strip}\label{sec:strip}
\subsection{Real exponential and hyperbolic functions}
We use the established Gonshor exponential
\[
 \exp_{\No}:(\No,+)\longrightarrow(\No_{>0},\cdot)
\]
and its inverse $\log_{\No}$ \cite{Gonshor,vdDE}. They are order-preserving group isomorphisms, agree with the ordinary real functions, and satisfy
\begin{equation}\label{eq:expreallocal}
 \exp_{\No}(x+h)=\exp_{\No}(x)\Ex(h)
 \qquad(h\in\mm_{\R}).
\end{equation}
No surreal field derivation is being chosen: derivatives below are derivatives of functions with respect to their argument.

For every $y\in\No$, define
\[
 \sinh y=\frac{\exp_{\No}y-\exp_{\No}(-y)}2,\qquad
 \cosh y=\frac{\exp_{\No}y+\exp_{\No}(-y)}2.
\]
The hyperbolic addition identities, $\cosh^2y-\sinh^2y=1$, and the usual derivatives follow directly from the group law and \eqref{eq:expreallocal}. In particular, $\sinh y=0$ exactly at $y=0$.

\subsection{Two complementary strips and an exponential quotient}
Define
\[
 \mathscr H=\{x+iy:x\in\No,\ y\in\OO_{\R}\},\qquad
 \mathscr V=\{x+iy:x\in\OO_{\R},\ y\in\No\}.
\]
These are additive subclasses of $\SC$, and $i\mathscr V=\mathscr H$. On the first strip define
\begin{equation}\label{eq:stripexp}
 E(x+iy)=\exp_{\No}(x)\cis y.
\end{equation}

\begin{theorem}[Canonical strip exponential]\label{thm:stripexp}
The map $E:\mathscr H\to\SC^\times$ is a surjective group homomorphism with kernel $2\pi i\Z$. It is fine-locally analytic and satisfies
\[
 E(z+h)=E(z)\Ex(h),\qquad E'(z)=E(z)
\]
for infinitesimal $h\in\SC$. Consequently
\begin{equation}\label{eq:stripquotient}
 \mathscr V/(2\pi\Z)\ \longrightarrow\ \SC^\times,
 \qquad [z]\longmapsto E(iz),
\end{equation}
is a group isomorphism.
\end{theorem}
\begin{proof}
The group law follows from the real exponential and finite-angle addition. If $E(x+iy)=1$, its norm forces $x=0$, and \cref{thm:polar} forces $y\in2\pi\Z$. For $w\ne0$, choose a finite angle $\theta$ of $w$; then
\[
 E(\log_{\No}|w|+i\theta)=w.
\]
The local formula follows by splitting an infinitesimal increment into its real and imaginary parts and using \eqref{eq:expreallocal}. The local power series gives the derivative. Multiplication by $i$ proves \eqref{eq:stripquotient}.
\end{proof}
The quotient is an algebraic statement, not a claim that this class strip is a classical topological universal covering space.

\begin{definition}[Canonical strip sine and cosine]
For $z\in\mathscr V$, set
\begin{equation}\label{eq:stripsincos}
 \Sin z=\frac{E(iz)-E(-iz)}{2i},\qquad
 \Cos z=\frac{E(iz)+E(-iz)}2.
\end{equation}
When $z$ is a finite real surreal, these are the lowercase functions already defined. When $z$ is finite surcomplex, these are the Taylor lifts of the ordinary entire complex sine and cosine.
\end{definition}

\begin{theorem}[Complex-variable identities and zero sets]\label{thm:striptrig}
For $z=x+iy\in\mathscr V$,
\begin{align}
 \Sin(x+iy)&=\sin x\cosh y+i\cos x\sinh y,\\
 \Cos(x+iy)&=\cos x\cosh y-i\sin x\sinh y,\\
 |\Sin(x+iy)|^2&=\sin^2x+\sinh^2y,\\
 |\Cos(x+iy)|^2&=\cos^2x+\sinh^2y.
 \label{eq:stripmoduli}
\end{align}
On this entire strip the addition identities, $\Sin^2z+\Cos^2z=1$, and
\[
 \Sin'=\Cos,\qquad\Cos'=-\Sin
\]
hold exactly. The zeros of $\Sin$ are precisely $\pi\Z$, and those of $\Cos$ are precisely $\pi/2+\pi\Z$; all are simple. The quotient $\Tan=\Sin/\Cos$ has simple poles at the latter points.
\end{theorem}
\begin{proof}
Insert \eqref{eq:stripexp} in \eqref{eq:stripsincos} and use the real hyperbolic definitions to obtain the first two formulas. Squaring real and imaginary parts and using $\cosh^2y=1+\sinh^2y$ gives the two norm identities. All other identities follow algebraically from the group law for $E$ and its local derivative.

A zero in either norm identity forces $\sinh y=0$, hence $y=0$. For finite real $x$, taking standard part and then the leading infinitesimal term shows that $\sin x=0$ exactly for $x\in\pi\Z$, and similarly for cosine. Equivalently, this follows from the kernel of $\cis$ and the double-angle identities. The derivatives at the zeros are $\pm1$, giving simplicity and the pole statement.
\end{proof}

\begin{theorem}[Surjectivity and complete fibers of strip sine]\label{thm:sinefibers}
Every $w\in\SC$ is $\Sin z$ for some $z\in\mathscr V$. Modulo $2\pi$, it has exactly two preimages counted with multiplicity; the two coincide exactly when $w=1$ or $w=-1$. For $z_1,z_2\in\mathscr V$,
\begin{equation}\label{eq:sinefibers}
 \Sin z_1=\Sin z_2
 \ \Longleftrightarrow\\
 z_1-z_2\in2\pi\Z\ \text{ or }\ z_1+z_2\in\pi+2\pi\Z.
\end{equation}
\end{theorem}
\begin{proof}
Put $u=E(iz)$. The equation $\Sin z=w$ becomes
\[
 u^2-2iwu-1=0.
\]
The algebraically closed field $\SC$ supplies two roots counted with multiplicity, both nonzero. Each has exactly one preimage modulo $2\pi$ by \cref{thm:stripexp}. The quadratic discriminant is $4(1-w^2)$, proving the multiplicity assertion. Alternatively, subtract the two sine values and use the identity
\[
 \Sin z_1-\Sin z_2
 =2\Cos\frac{z_1+z_2}{2}\,\Sin\frac{z_1-z_2}{2}.
\]
The zero sets in \cref{thm:striptrig} give \eqref{eq:sinefibers}.
\end{proof}
Thus canonical complex trigonometry is not confined to finite surcomplex values or finite imaginary arguments. Its unresolved phase choice is specifically at \emph{infinite real arguments}.

\section{Trigonometric polynomials and sharp algebraic root counts}\label{sec:trigpoly}
Let $n\in\N$ be ordinary and let $c_{-n},\ldots,c_n\in\SC$. Define on $\mathscr V$
\begin{equation}\label{eq:trigpoly}
 T(z)=\sum_{k=-n}^{n}c_kE(ikz).
\end{equation}
This is a finite sum, so arbitrary surreal sizes of coefficients cause no summability problem. On finite real angles it is a trigonometric polynomial. It is real-valued there exactly when $c_{-k}=\bar c_k$ for all $k$; the converse follows by the polynomial identity principle below.

\begin{theorem}[Algebraization and the $2n$ root bound]\label{thm:polyroots}
Suppose $T\ne0$ and set
\begin{equation}\label{eq:algebraization}
 P(U)=\sum_{k=-n}^{n}c_kU^{k+n}.
\end{equation}
Under $U=E(iz)$, zeros of $T$ modulo $2\pi$ correspond bijectively, with multiplicity, to nonzero roots of $P$. Hence $T$ has at most $2n$ zeros modulo $2\pi$, counted with multiplicity, on all of $\mathscr V$. If $c_{-n}c_n\ne0$, the count is exactly $2n$. In particular, the number of finite real angle zeros is at most $2n$.
\end{theorem}
\begin{proof}
The identity $T(z)=U^{-n}P(U)$ and the bijection \eqref{eq:stripquotient} prove the set-theoretic correspondence. Near a point with value $U_0\ne0$, the map is
\[
 E(i(z+h))=U_0\Ex(ih)=U_0+iU_0h+O(h^2)
\]
with nonzero linear coefficient. Substitution by such a formal local coordinate preserves the order of a zero. Polynomial factorization in $\SC$ gives the total number of nonzero roots as $\deg P-\operatorname{ord}_0P\le2n$. Both endpoint coefficients being nonzero makes this number $2n$.
\end{proof}
There cannot be a nonzero finite trigonometric polynomial vanishing at every ordinary real angle: its associated polynomial would have infinitely many distinct ordinary unit-circle roots. This observation also proves uniqueness of the finite Fourier coefficients over $\SC$.

\begin{corollary}[Chebyshev equations and extrema equations]
For an ordinary $n\ge1$ and $w\in\SC$, the equation $\Cos(nz)=w$ has exactly $2n$ solutions modulo $2\pi$, counted with multiplicity. For a nonconstant real trigonometric polynomial of degree at most $n$, its derivative has at most $2n$ finite real zeros modulo $2\pi$.
\end{corollary}
\begin{proof}
For the first assertion, the associated equation is $U^{2n}-2wU^n+1=0$, with nonzero constant and leading coefficients. For the second, differentiation multiplies coefficient $c_k$ by $ik$, preserves the band, and does not annihilate a nonconstant polynomial in characteristic zero. Apply \cref{thm:polyroots}.
\end{proof}
The derivative statement bounds stationary angles; it does not invoke a general compactness or extreme-value theorem on a surreal interval.

\section{Support-controlled lifting of trigonometric zero clusters}\label{sec:clusters}
This section specializes the positive-support philosophy of the supplied finite-intersection articles. The proof is polynomial and is included to avoid making the geometric theory depend on their more substantial analytic division results.

Let
\[
 T_0(z)=\sum_{k=-n}^{n}a_kE(ikz),\qquad a_k\in\C,
 \qquad a_{-n}a_n\ne0,
\]
where on ordinary complex arguments this is the usual trigonometric polynomial. Perturb the coefficients by $e_k\in\mm_{\C}$, allowing zero values:
\[
 T(z)=\sum_{k=-n}^{n}(a_k+e_k)E(ikz).
\]
Let $S\subset\Gamma_{>0}$ be the union of the supports of the nonzero $e_k$.

\begin{theorem}[Exact cluster multiplicity and coefficient support]\label{thm:cluster}
Suppose $\theta_0\in\C$ is an ordinary zero of $T_0$ of multiplicity $m$, and put $u_0=e^{i\theta_0}$. Then $T$ has exactly $m$ zeros, counted with multiplicity, in the complex infinitesimal monad
\[
 \theta_0+\mm_{\C}.
\]
Normalize the algebraized polynomial to be monic. Its unique monic factor $A$ collecting the roots in $u_0+\mm_{\C}$ has degree $m$ and reduces to $(U-u_0)^m$. The coefficients of $A-(U-u_0)^m$, and of the complementary factor minus its ordinary reduction, are supported in $S^*\setminus\{0\}$.

If $\theta_0$ is real, both $T_0$ and $T$ are real on finite real arguments, and $m=1$, the unique lifted zero is real. No corresponding reality conclusion is asserted for an arbitrary multiple zero.
\end{theorem}
\begin{proof}
Divide $P(U)=\sum(a_k+e_k)U^{k+n}$ by $a_n+e_n$. Its inverse leading coefficient is an ordinary nonzero scalar times a geometric series in an infinitesimal; the normalized perturbation therefore has support in $S^*$. Its ordinary reduction has a coprime monic factorization
\[
 P_0=A_0B_0,\qquad A_0=(U-u_0)^m,\qquad B_0(u_0)\ne0.
\]
Seek monic factors $A=A_0+\Delta A$, $B=B_0+\Delta B$ of the same degrees, with positive-support coefficient corrections. The finite-dimensional ordinary linear map
\begin{equation}\label{eq:sylvester}
 (X,Y)\longmapsto XB_0+A_0Y,
 \qquad \deg X<m,\quad \deg Y<2n-m,
\end{equation}
is an isomorphism onto the polynomials of degree below $2n$. Indeed, its kernel is zero by coprimality and the degree restrictions, and the dimensions agree.

At a positive exponent $\nu\in S^*$, the coefficient of $AB=P$ is a linear equation with map \eqref{eq:sylvester}; its right side is the prescribed coefficient of $P-P_0$ minus the already determined coefficient of $\Delta A\Delta B$. In that latter product, the two exponents are positive and strictly below $\nu$. Transfinite recursion on the well-ordered monoid $S^*$ therefore constructs the factors. Each equation involves only finitely many terms by \cref{lem:support}. The same recursion proves uniqueness. This is a coefficient recursion, not a convergence claim for a Newton sequence.

Translate $U=u_0+H$. The polynomial $A(u_0+H)$ is monic of degree $m$ with all lower coefficients infinitesimal. Every one of its roots is infinitesimal: if a root $h$ were appreciable or infinite, division of its equation by $h^m$ would give $1$ plus a finite sum of infinitesimals equal to zero. Its $m$ roots exist in an algebraically closed divisible Hahn workspace. The complementary factor cannot vanish at $u_0+h$ for infinitesimal $h$, since its standard part there is $B_0(u_0)\ne0$.

The local bijection
\[
 U\longmapsto \theta_0-i\Logm(U/u_0)
\]
transports these roots to the specified angle monad. Its nonzero linear term preserves multiplicity. Finally, conjugation preserves a real trigonometric equation and the monad of a real $\theta_0$. When there is exactly one root, that root must be fixed by conjugation.
\end{proof}

The support bound belongs to \emph{factor coefficients}, not to individual roots. A root may require a fractional scale such as $t^{\lambda/2}$ even when the coefficient perturbation only uses $t^\lambda$. This is exactly what tangency in \cref{thm:tangency} exhibits.

\begin{example}[A double zero splitting off the real angle axis]
The equation $1-\Cos z=-\eps$, with $\eps>0$ infinitesimal, has two zeros in the monad of zero. They are
\[
 z=\pm i\,\arcosh(1+\eps),
 \qquad
 \arcosh(1+\eps)=\log_{\No}\left(1+\eps+\sqrt{2\eps+\eps^2}\right).
\]
They have valuation $\tfrac12v(\eps)$ and are not real. Indeed, $\Cos(iy)=\cosh y$ and the defining inverse hyperbolic formula gives $\cosh y=1+\eps$. This is the complex continuation of the disappearance of real intersections after tangency.
\end{example}

\section{Finite Fourier calculus and positivity on the surreal circle}\label{sec:fourier}
\subsection{Coefficientwise integration and Parseval}
For a finite trigonometric polynomial, define its Hahn integral by expanding its finitely many coefficients in normal form and integrating each ordinary coefficient function over an ordinary period. This is legitimate because the support is a finite union of well-ordered coefficient supports. In particular,
\begin{equation}\label{eq:orthogonality}
 \frac1{2\pi}\HInt_0^{2\pi}E(ik\theta)\,d\theta
 =\begin{cases}1&k=0,\\0&k\ne0,\end{cases}
 \qquad k\in\Z.
\end{equation}
Here the ordinary coefficient integral is the usual real or complex integral; its Hahn assembly is not an ordinary Riemann limit over the full class interval.

\begin{theorem}[Finite Fourier identities over $\SC$]\label{thm:parseval}
If $F(\theta)=\sum_{k=-n}^nf_k\cis(k\theta)$ and
$G(\theta)=\sum_{k=-m}^mg_k\cis(k\theta)$, extending coefficients by zero when needed, then
\begin{align}
 f_j&=\frac1{2\pi}\HInt_0^{2\pi}F(\theta)\cis(-j\theta)\,d\theta,\\
 \frac1{2\pi}\HInt_0^{2\pi}F(\theta)\overline{G(\theta)}\,d\theta
 &=\sum_k f_k\bar g_k,\\
 \frac1{2\pi}\HInt_0^{2\pi}|F(\theta)|^2\,d\theta
 &=\sum_k|f_k|^2.
 \label{eq:parseval}
\end{align}
In particular, the last expression is positive unless $F=0$.
\end{theorem}
\begin{proof}
Expand the finite sums, use $\overline{\cis(k\theta)}=\cis(-k\theta)$, and apply \eqref{eq:orthogonality}. Every operation on the frequency indices is finite. Positivity follows from the ordering of $\No$ and positivity of each squared norm.
\end{proof}
These identities remain valid for infinite coefficients and for infinitesimal coefficients. They do not establish convergence or completeness for an unrestricted infinite Fourier series. In particular, infinitely many ordinary Fourier summands at the same Hahn exponent do not form a strongly summable family merely because the corresponding real Fourier series converges.

\subsection{Fej\'er--Riesz factorization with surreal coefficients}
The classical Fej\'er--Riesz theorem factors a nonnegative trigonometric polynomial as the squared modulus of a polynomial on the circle; see \cite{DR} for its established context. In the present scalar setting a proof over the real closed surreal field is particularly direct.

\begin{lemma}[Nonnegative polynomials over a real closed field]\label{lem:twosquares}
Let $F$ be real closed, let $K=F[i]$, and let $p\in F[t]$ be nonnegative at every $t\in F$. Then there is $q\in K[t]$ such that
\[
 p(t)=q(t)\overline{q(t)}\quad(t\in F).
\]
If $p\ne0$, then $\deg q=\tfrac12\deg p$.
\end{lemma}
\begin{proof}
For $p=0$ take $q=0$. Otherwise factor $p$ over $F$ into real linear factors and irreducible quadratic factors. Every real root has even multiplicity: an odd multiplicity would change the sign across that root, as is seen by taking a sufficiently small increment avoiding the finitely many other roots. The leading constant is positive, and each monic irreducible quadratic has form $(t-a)^2+b^2$ with $b>0$. Take one factor $t-a-ib$ from each conjugate pair, half the multiplicity of each real linear factor, and the positive square root of the leading constant. Their product is the required $q$.
\end{proof}

\begin{theorem}[Surreal Fej\'er--Riesz theorem]\label{thm:fejer}
Let
\[
 T(\theta)=\sum_{k=-n}^{n}c_k\cis(k\theta),\qquad
 c_{-k}=\bar c_k,
\]
with coefficients in $\SC$. The following are equivalent:
\begin{enumerate}
 \item $T(\theta)\ge0$ for every finite surreal angle $\theta$.
 \item There is a polynomial $Q(U)=\sum_{j=0}^nq_jU^j\in\SC[U]$ such that
 \begin{equation}\label{eq:fejer}
 T(\theta)=|Q(\cis\theta)|^2
 \qquad\text{for every finite surreal }\theta.
 \end{equation}
\end{enumerate}
\end{theorem}
\begin{proof}
The second condition plainly implies the first. Conversely use the rational parametrization $u(t)$ in \eqref{eq:cayley} and form
\begin{equation}\label{eq:cayleypositive}
 p(t)=(1+t^2)^n\sum_{k=-n}^{n}c_ku(t)^k.
\end{equation}
Since $u(t)=(1+it)^2/(1+t^2)$ and $u(t)^{-1}=(1-it)^2/(1+t^2)$, the denominators cancel and $p$ is a polynomial of degree at most $2n$. Its values and coefficients are real, by the Hermitian symmetry of the $c_k$. The hypothesis and the surjectivity of the rational parametrization give $p(t)\ge0$ for every $t\in\No$.

By \cref{lem:twosquares}, $p(t)=|q(t)|^2$ with $q\in\SC[t]$ of degree at most $n$. Define
\begin{equation}\label{eq:fejerconstruction}
 Q(U)=\left(\frac{1+U}{2}\right)^n
 q\left(\frac{U-1}{i(U+1)}\right).
\end{equation}
The degree bound on $q$ makes $Q$ a polynomial of degree at most $n$. For $U=u(t)$ one has $(1+U)/2=1/(1-it)$; hence
\[
 |Q(u(t))|^2=\frac{|q(t)|^2}{(1+t^2)^n}=T(u(t)),
\]
where $T(U)$ denotes the Laurent polynomial associated with $T(\theta)$. This proves equality on every unit point except possibly $U=-1$. On the circle, replace conjugate powers by negative powers. The difference is then a Laurent polynomial; after multiplication by $U^n$ it is an ordinary-degree polynomial with infinitely many roots, so it vanishes identically. Equality holds at $U=-1$ as well.
\end{proof}
The proof is finite algebra over a real closed field. It uses no real compactness theorem, no Hilbert-space completion, and no limit of spectral factors.

\begin{corollary}[Fourier coefficient bounds from positivity]\label{cor:coeffbounds}
Under the nonnegativity hypothesis of \cref{thm:fejer},
\[
 c_0=\sum_{j=0}^n|q_j|^2,\qquad |c_k|\le c_0\quad(1\le k\le n).
\]
If $n\ge1$, the highest coefficient satisfies the sharper bound $2|c_n|\le c_0$.
\end{corollary}
\begin{proof}
Expanding \eqref{eq:fejer} gives
\[
 c_k=\sum_{j=0}^{n-k}q_{j+k}\bar q_j\qquad(k\ge0).
\]
Finite-dimensional Cauchy--Schwarz, proved over $\No$ by expanding a sum of squared norms, gives
\[
 |c_k|^2\le\left(\sum_{j=0}^{n-k}|q_{j+k}|^2\right)
                \left(\sum_{j=0}^{n-k}|q_j|^2\right)\le c_0^2.
\]
For $k=n$, $c_n=q_n\bar q_0$ and
$2|q_n||q_0|\le|q_n|^2+|q_0|^2\le c_0$.
\end{proof}

\begin{example}[Ordinary-angle testing misses infinitesimal negativity]\label{ex:positivity}
Let $0<\eps\in\mm_{\R}$ and consider the degree-two polynomial
\[
 T(\theta)=(\cos\theta-\eps)^2-\eps^4.
\]
For every \emph{ordinary real} $\theta$, $T(\theta)>0$. If $\cos\theta\ne0$, its standard part is the positive real $\cos^2\theta$; if $\cos\theta=0$, its value is $\eps^2-\eps^4>0$. But at the finite nonordinary angle $\theta=\arccos\eps$,
\[
 T(\theta)=-\eps^4<0.
\]
Thus testing only the ordinary circle is not sufficient for the hypothesis of \cref{thm:fejer}, even for a fixed, very small trigonometric degree. A squared-modulus factorization valid on the full surreal circle cannot exist for this example.
\end{example}

\section{Hyperbolic trigonometry in the full surreal disk}\label{sec:hyperbolic}
\subsection{A disk metric that includes infinitesimal boundary distances}
Let
\[
 \mathbb D_{\No}=\{z\in\SC:|z|<1\}.
\]
Unlike the ordinary-disk thickening used in some coherent analytic statements, this disk includes points infinitesimally inside the unit circle whose standard parts lie on its boundary. For $a\in\mathbb D_{\No}$ define
\begin{equation}\label{eq:diskmap}
 \varphi_a(z)=\frac{z-a}{1-\bar a z}.
\end{equation}
The denominator is nonzero because $|\bar a z|<1$. Direct expansion gives
\begin{equation}\label{eq:diskidentity}
 1-|\varphi_a(z)|^2
 =\frac{(1-|a|^2)(1-|z|^2)}{|1-\bar a z|^2}>0.
\end{equation}
Thus $\varphi_a$ is a disk automorphism, with inverse $(w+a)/(1+\bar a w)$. Rotations are disk automorphisms as well.

Define
\begin{align}
 \rho(z,w)&=\left|\frac{z-w}{1-\bar w z}\right|,\\
 d_H(z,w)&=\log_{\No}\frac{1+\rho(z,w)}{1-\rho(z,w)}
 =2\artanh\rho(z,w),
 \label{eq:diskdistance}
\end{align}
where $\artanh r=\tfrac12\log_{\No}((1+r)/(1-r))$ for $|r|<1$.

\begin{proposition}[Disk invariance and metric properties]\label{prop:diskmetric}
The quantities $\rho$ and $d_H$ are invariant under rotations and the maps $\varphi_a$. The function $d_H$ is a $\No$-valued metric on $\mathbb D_{\No}$.
\end{proposition}
\begin{proof}
The invariance of $\rho$ is obtained by substituting \eqref{eq:diskmap}, putting differences over common denominators, and canceling absolute values. Invariance of $d_H$ follows. Symmetry, positivity, and separation of points follow from the definition and \eqref{eq:diskidentity}.

For the triangle inequality, move the middle point to zero. If $|z|=r$, $|w|=s$, and the angle between them is $\theta$, then
\[
 \rho(z,w)^2=\frac{r^2+s^2-2rs\cos\theta}
                   {1+r^2s^2-2rs\cos\theta}
 \le\left(\frac{r+s}{1+rs}\right)^2.
\]
For fixed $r,s$, the ratio on the left is largest when $\cos\theta=-1$: write it as $1-(1-r^2)(1-s^2)/(1+r^2s^2-2rs\cos\theta)$ and compare denominators. Applying the increasing logarithm and using
\[
 \frac{1+(r+s)/(1+rs)}{1-(r+s)/(1+rs)}
 =\frac{1+r}{1-r}\frac{1+s}{1-s}
\]
gives $d_H(z,w)\le d_H(z,0)+d_H(0,w)$.
\end{proof}

At a vertex $A$, define the hyperbolic angle between the sides towards $B,C$ to be the Euclidean angle between $\varphi_A(B)$ and $\varphi_A(C)$. This is the angle between the corresponding circular arcs orthogonal to the boundary, or diameters when appropriate. The definition is invariant under disk automorphisms: after moving the vertex to zero, the remaining automorphism fixing zero is a rotation, as direct algebra with the fractional-linear formulas shows. Equivalently, the nonzero complex derivative
\[
 \varphi_a'(z)=\frac{1-|a|^2}{(1-\bar a z)^2}
\]
preserves angles. All these angles are finite, even when their side distances are infinite.

\subsection{Hyperbolic cosine and sine laws}
\begin{theorem}[Hyperbolic triangle laws over $\No$]\label{thm:hyperbolic}
Let $A,B,C$ be distinct disk points not on one hyperbolic line. Write
\[
 a=d_H(B,C),\quad b=d_H(C,A),\quad c=d_H(A,B),
\]
and let $\alpha,\beta,\gamma\in(0,\pi)$ be their hyperbolic interior angles. Then
\begin{align}
 \cosh a&=\cosh b\cosh c-\sinh b\sinh c\cos\alpha,
 \label{eq:hypcos}\\
 \frac{\sin\alpha}{\sinh a}
 &=\frac{\sin\beta}{\sinh b}
 =\frac{\sin\gamma}{\sinh c}.
 \label{eq:hypsine}
\end{align}
The distances $a,b,c$ may be infinite.
\end{theorem}
\begin{proof}
Move $A$ to zero and rotate $B$ to $r>0$. Write $C=s\cis\alpha$ after a possible orientation reversal. Here $r=\tanh(c/2)$ and $s=\tanh(b/2)$, so
\[
 \cosh c=\frac{1+r^2}{1-r^2},\quad
 \sinh c=\frac{2r}{1-r^2},
\]
and similarly for $b$. Put $q=\rho(B,C)$. Then
\[
 q^2=\frac{r^2+s^2-2rs\cos\alpha}
           {1+r^2s^2-2rs\cos\alpha},\qquad
 \cosh a=\frac{1+q^2}{1-q^2}.
\]
Substitution yields
\[
 \cosh a=
 \frac{(1+r^2)(1+s^2)-4rs\cos\alpha}
 {(1-r^2)(1-s^2)},
\]
which is \eqref{eq:hypcos}. By invariance the same calculation applies at every vertex.

For the sine law, set $X=\cosh a$, $Y=\cosh b$, $Z=\cosh c$. The cosine law gives
\[
 \sin^2\alpha
 =\frac{1+2XYZ-X^2-Y^2-Z^2}{(Y^2-1)(Z^2-1)}.
\]
Divide by $\sinh^2a=X^2-1$. The result is symmetric in $a,b,c$, so the squares of all three ratios in \eqref{eq:hypsine} agree. The ratios are strictly positive, giving the unsquared equality.
\end{proof}

\begin{example}[An infinitely distant point inside the disk]
For $0<\eps\in\mm_{\R}$ and $z=1-\eps$,
\[
 d_H(0,z)=\log_{\No}\frac{2-\eps}{\eps}
 =-\log_{\No}\eps+\log 2+\log_{\No}(1-\eps/2).
\]
This distance is infinite. The last logarithm is a strongly summable infinitesimal series. No value of sine or cosine at an infinite \emph{real circular angle} is required to obtain either this distance or the triangle laws.
\end{example}

\section{Coherent arc length, sector area, and the limit pitfall}\label{sec:arcs}
\subsection{An explicit integral convention}
For an ordinary analytic $f$ with a primitive $F$ on a real interval and finite endpoints $a=a_0+\eps$, $b=b_0+\eta$, define its lifted integral by
\begin{align}
 \HInt_a^b f\sh(\theta)\,d\theta
 &=F\sh(b)-F\sh(a)\label{eq:endpointintegral}\\
 &=\int_{a_0}^{b_0}f(t)\,dt
 +\sum_{m\ge1}\frac{F^{(m)}(b_0)}{m!}\eta^m
 -\sum_{m\ge1}\frac{F^{(m)}(a_0)}{m!}\eps^m.\notag
\end{align}
The constant of the primitive cancels. More general common-domain Hahn families are integrated coefficientwise, with the same endpoint corrections whenever the supports justify evaluation. This is the convention needed below and is consistent with the supplied analysis manuscript. It is not a definition of an unrestricted Riemann integral on all surreal functions. Other surreal integration frameworks have substantially different scopes and hypotheses \cite{CE}.

\begin{theorem}[Circle length and sector area at every radius]\label{thm:arcs}
For $O\in\SC$, $R\in\No_{>0}$, and finite $a\le b$, consider the parametrized arc
\[
 z(\theta)=O+R\cis\theta,\qquad a\le\theta\le b.
\]
Its coherent arc length, computed from its speed with \eqref{eq:endpointintegral}, is $R(b-a)$. Its oriented sector area about $O$ is $R^2(b-a)/2$. In particular, the circle has coherent circumference $2\pi R$ and disk area $\pi R^2$.
\end{theorem}
\begin{proof}
One has $z'(\theta)=iR\cis\theta$, so $|z'|=R$. The integral of this constant is $R(b-a)$, including infinitesimal endpoint corrections. The usual parametrized sector functional has integrand
\[
 \tfrac12\im\bigl(\overline{z-O}\,z'\bigr)=\frac{R^2}{2},
\]
whose integral is $R^2(b-a)/2$. All expressions are finite algebraic combinations of a fixed surreal scalar and scalar-lifted trigonometric functions, so their coefficientwise integrals have valid common support. Taking $b-a=2\pi$ proves the circle statements.
\end{proof}
As a related contour identity, direct parametrization gives
\[
 \HInt_{|z-O|=R}\frac{dz}{z-O}=2\pi i.
\]
This is a coherent contour functional, not an assertion about ordinary paths being connected in the full fine topology.

\subsection{Why ordinary polygon limits do not define this length}
\begin{proposition}[Failure of fine convergence of inscribed perimeters]\label{prop:perimeters}
Let $R>0$ be any surreal radius and let
\[
 P_n=2nR\sin(\pi/n),\qquad n=3,4,\ldots,
\]
be the perimeters of ordinary finite regular inscribed polygons. The sequence does not converge to $2\pi R$ in the full surreal fine topology. Moreover, the set $\{P_n:n\ge3\}$ has no least upper bound in $\No$.
\end{proposition}
\begin{proof}
The ordinary real numbers $q_n=P_n/R=2n\sin(\pi/n)$ satisfy $q_n<2\pi$ and converge to $2\pi$ in $\R$. For any positive surreal infinitesimal $\eta$, the positive ordinary gap $2\pi-q_n$ is greater than $\eta$ for every $n$. Hence
\[
 2\pi R-P_n>R\eta
\]
for every $n$, ruling out fine convergence.

An infinite upper bound for the $q_n$ cannot be least because $2\pi$ is a smaller upper bound. A finite candidate with standard part below $2\pi$ is not an upper bound by ordinary real convergence. A candidate with standard part above $2\pi$ can be decreased. Finally, every finite candidate with standard part exactly $2\pi$ can be decreased by a positive infinitesimal and remain an upper bound: each ordinary gap is still larger than any such correction. There is therefore no least upper bound. Multiplication by $R$ is an order isomorphism and preserves this conclusion.
\end{proof}
Thus the circumference formula in \cref{thm:arcs} is a genuine coherent-integration theorem. It must not be justified by importing a real supremum or a real sequence limit into $\No$ without specifying the interpretation.

\section{Infinite real phases: extensions, freedom, and obstruction}\label{sec:global}
\subsection{Classification of circular group-law extensions}
Split every surreal normal form into its negative-exponent and nonnegative-exponent parts:
\begin{equation}\label{eq:splitinfinite}
 x=p+\fin(x),\qquad p\in\Pi,\quad\fin(x)\in\OO_{\R},
\end{equation}
where $\Pi$ is the additive class of purely infinite surreals. This is an additive direct-sum decomposition $\No=\Pi\oplus\OO_{\R}$. The finite part retains both the ordinary constant and all infinitesimal terms; it is not merely standard part.

\begin{theorem}[All circular group-law extensions]\label{thm:characters}
For every group homomorphism $\chi:(\Pi,+)\to(\TT(\No),\cdot)$, the formula
\begin{equation}\label{eq:globalcis}
 U_\chi(p+u)=\chi(p)\cis u,\qquad u\in\OO_{\R},
\end{equation}
defines a group homomorphism $\No\to\TT(\No)$ extending canonical finite-angle $\cis$. Every such extension is uniquely of this form.

Writing $U_\chi=C_\chi+iS_\chi$, one obtains real-valued global sine and cosine satisfying the addition formulas, $C_\chi^2+S_\chi^2=1$, and, in fine-local calculus,
\[
 S_\chi'=C_\chi,\qquad C_\chi'=-S_\chi.
\]
For every infinitesimal real $h$,
\begin{equation}\label{eq:globallocal}
 U_\chi(x+h)=U_\chi(x)\Ex(ih).
\end{equation}
\end{theorem}
\begin{proof}
The direct-sum decomposition and the two group laws prove that \eqref{eq:globalcis} is a homomorphism with the desired restriction. Conversely, the restriction of any extension to $\Pi$ determines $\chi$, and its value at $p+u$ must be the product of its two restrictions. Formula~\eqref{eq:globallocal} follows because an infinitesimal increment belongs to the finite summand. Expanding its right side proves fine-local analyticity and the derivative identities.
\end{proof}

\begin{example}[Any prescribed phase at $\omega$]\label{ex:phases}
Let $\ell(p)\in\R$ be the coefficient of $\omega=t^{-1}$ in a purely infinite normal form. For any finite angle $\tau\in\OO_{\R}$,
\[
 \chi_\tau(p)=\cis\bigl(\tau\ell(p)\bigr)
\]
is a character, because $\ell$ is additive and $\tau\ell(p)$ is finite. Its extension satisfies $U_{\chi_\tau}(\omega)=\cis\tau$. Thus $S(\omega)$ and $C(\omega)$ can jointly be assigned any point of the surreal unit circle without changing the finite-angle theory, addition laws, or fine-local differential equations. For example, phases $0$, $\pi/2$, and $\pi$ give respectively $(C(\omega),S(\omega))=(1,0),(0,1),(-1,0)$.
\end{example}
The supplied analysis manuscript already constructs an ordinary-real subfamily of this kind for global complex exponentials. The present classification exposes the entire character freedom relevant to circular trigonometry. The character $\chi=1$ gives the explicit normal-form extension $U_0(x)=\cis(\fin(x))$. It is a possible convention, not a consequence of the local trigonometric axioms alone.

\begin{theorem}[Infinite periods are unavoidable]\label{thm:infiniteperiods}
Every homomorphism extension in \cref{thm:characters} has a kernel element in every coset $p+\OO_{\R}$ of the finite subgroup. In particular, for every nonzero $p\in\Pi$ it has an infinite period. No extension agreeing with canonical finite-angle $\cis$ can have kernel exactly $2\pi\Z$.
\end{theorem}
\begin{proof}
Finite-angle $\cis$ is already surjective onto the whole unit circle. Choose finite $a$ with $\cis a=\chi(p)^{-1}$. Then $U_\chi(p+a)=1$. For nonzero purely infinite $p$, $p+a$ is infinite. Being in the kernel of a group homomorphism makes it a period of both real and imaginary components.
\end{proof}
This is stronger than a bare example of nonuniqueness. The ordinary statement ``all periods are ordinary multiples of $2\pi$'' is incompatible with extending the circular group law to all surreal real inputs while retaining its finite-angle meaning.

Each $U_\chi$ also gives a global complex exponential
\[
 E_\chi(x+iy)=\exp_{\No}(x)U_\chi(y),
\]
and corresponding global complex sine and cosine by the usual exponential formulas. All these agree with the canonical strip functions of \cref{sec:strip}. Their disagreement occurs beyond that strip.

\subsection{Why infinite-frequency scaling destroys common-domain coherence}
A common-domain Hahn-analytic function on an ordinary neighborhood $J$ of zero is a family
\[
 H=\sum_{\gamma\in S}h_\gamma t^\gamma
\]
with one well-ordered support and ordinary real-analytic coefficients on the same $J$; evaluation at finite thickened points uses Taylor substitution. One may equally use ordinary holomorphic coefficients on a complex neighborhood. This is the real restriction of the coherence convention in the supplied manuscripts.

\begin{theorem}[No coherent infinite-frequency bounded sine]\label{thm:infinitefrequency}
Let $S_\chi$ be any global extension from \cref{thm:characters}, and let $L>0$ be infinite. There is no common-domain Hahn-analytic $H$ on an ordinary neighborhood of zero whose evaluation satisfies
\[
 H\sh(x)=S_\chi(Lx)
\]
throughout its finite surreal thickening.
\end{theorem}
\begin{proof}
For every ordinary real $r$ in the neighborhood, the right side belongs to $[-1,1]$. Therefore every negative-exponent coefficient $h_\gamma(r)$ of $H(r)$ is zero, by uniqueness of normal form and finiteness. This holds for all such $r$, so every negative-exponent coefficient function vanishes identically. The derivative $H'(0)=\sum h_\gamma'(0)t^\gamma$ consequently has no negative exponents and is finite.

But equality on the thickening implies equality of the fine derivatives at zero. From \eqref{eq:globallocal}, or from the chain rule with the constant scalar $L$,
\[
 \left.\frac{d}{dx}S_\chi(Lx)\right|_{x=0}=L,
\]
which is infinite. This contradiction proves the theorem.
\end{proof}
The requirement that the coefficient functions share an \emph{ordinary} neighborhood is essential. On a sufficiently small surreal scale, every $S_\chi(Lx)$ still has its usual local power series. The failure is the compatibility of that infinite-frequency behavior with a bounded, common-domain ordinary analytic coefficient description.

\section{Conclusions, dependencies, and computational verification}\label{sec:conclusion}
\subsection{What the theory establishes}
The central geometric conclusion is exact: arbitrary surreal lengths do not require infinite circular angles. The finite-angle quotient describes every unit direction, and ordinary finite triangle constructions give the same identities at all scales. The infinitesimal part is not negligible data; it controls chord precision, nearly straight angles, the order of the triangle-inequality slack, and tangency splitting.

A second conclusion concerns the complex plane. The canonical strip $\mathscr V$ is large enough for a surjective complex sine and for the complete algebraic solution theory of every finite trigonometric polynomial. The passage $U=E(iz)$ turns periodic root questions into finite-degree polynomial questions. Positive-support coefficient lifting distinguishes stable finite algebras from unstable individual roots. Positivity also has to be tested on the full surreal circle, not only on its ordinary real shadow.

Finally, full real-argument oscillation is a different problem. It admits many fine-analytic solutions, parameterized by characters of the purely infinite additive summand. These extensions necessarily introduce infinite periods and fail common-domain coherence at infinite frequency. Thus neither nonexistence nor canonical uniqueness is an appropriate blanket claim about ``surreal sine.'' The answer depends on the domain and the analytic structure specified.

\subsection{A dependency audit}
\begin{center}
\small
\begin{tabular}{L{0.28\textwidth}L{0.62\textwidth}}
\toprule
Result family & Inputs actually used\\
\midrule
Finite sine, cosine, and circle logarithm & Ordinary analytic germs, normal forms, positive-support Neumann lemma\\
Triangle and cyclic geometry & Ordered-field algebra, positive square roots, finite-angle circle parametrization\\
Valuation estimates and tangency & Exact algebraic identities and unit-leading infinitesimal series\\
Strip functions and hyperbolic disk & Gonshor exponential and logarithm, plus the finite-angle theory\\
Trigonometric root bounds & Algebraic closedness and the strip exponential quotient\\
Cluster lifting & Coprime polynomial factorization, a finite Sylvester map, transfinite coefficient recursion\\
Fej\'er--Riesz factorization & Real closedness, rational circle parametrization, finite polynomial factorization\\
Arc and sector formulas & Explicit coefficientwise/primitive-based integral convention\\
Infinite phase classification & The additive normal-form split and circle surjectivity\\
\bottomrule
\end{tabular}
\end{center}
No result uses a proper-class sum, assumes that a positive valuation increment is cofinal, or identifies fine differentiation with a chosen derivation of the surreal field. The Berarducci--Mantova differential-field structure \cite{BM} is important background in surreal analysis but is not an input to these function-derivative arguments.

\subsection{Reproducible checks and their limits}
The accompanying Python program \texttt{verify\_examples.py} uses exact SymPy algebra and formal power series. It checks the Gram, rational-circle, Heron, circumradius, Euler, and Ptolemy identities; Chebyshev formulas; flat-triangle and inverse-cosine coefficients; the algebraization of trigonometric equations; Fourier convolution and a spectral-factor example; and the rational identities behind the hyperbolic laws. All 68 checks passed in the accompanying run; the report lists each check individually.

These checks validate the finite algebra and selected Taylor coefficients occurring in the proofs. They do not implement the full surreal class, verify transfinite support arguments, establish the real closedness theorem, or machine-check the mathematical proofs. The support lemma, branch arguments, and proper-class distinctions are justified in the text, not inferred from numerical experiments.

\appendix
\section{A compact formula sheet with domain reminders}\label{app:formulas}
For finite real $\theta$, nonzero $z,w\in\SC$, and an ordinary integer $n$:
\begin{align*}
 z&=|z|\cis\Arg z &&\text{with }\Arg z\in\OO_{\R}/2\pi\Z,\\
 \cos\ang zw&=\frac{\re(\bar zw)}{|z||w|},&
 \sin\ang zw&=\frac{\im(\bar zw)}{|z||w|},\\
 \cis(n\theta)&=(\cis\theta)^n,&
 \tan(\theta/2)&=\frac{\sin\theta}{1+\cos\theta},\\
 |u-v|&=2\sin\bigl(d_\theta(u,v)/2\bigr)&& (u,v\in\TT(\No)),\\
 \arctan t&=\Arg(1+it)\in(-\pi/2,\pi/2)&& (t\in\No).
\end{align*}
The first argument is an angle class; the sine in the directed-angle formula may be negative. In a triangle with positive area, interior sines are positive and
\begin{gather*}
 \alpha+\beta+\gamma=\pi,\qquad
 a^2=b^2+c^2-2bc\cos\alpha,\qquad
 a=2R\sin\alpha,\\
 \Delta^2=s(s-a)(s-b)(s-c),\qquad
 r=\Delta/s,\qquad R=abc/(4\Delta),\\
 \tan(\alpha/2)=r/(s-a),\qquad |O-I|^2=R(R-2r).
\end{gather*}
For $z=x+iy$ with finite $x$ and arbitrary $y\in\No$,
\[
 \Sin z=\sin x\cosh y+i\cos x\sinh y,
 \qquad E(iz)=\exp_{\No}(-y)\cis x.
\]
For infinite real $x$, circular sine and cosine require an explicitly selected extension; the lowercase finite-angle notation alone does not specify their values.

\section{Worked formal series}\label{app:series}
All the following are strongly summable for infinitesimal $h$ or $\delta$, with the indicated positive square root in the last line:
\begin{align*}
 \arctan h&=h-\frac{h^3}{3}+\frac{h^5}{5}-\frac{h^7}{7}+O(h^9),\\
 \arcsin h&=h+\frac{h^3}{6}+\frac{3h^5}{40}+\frac{5h^7}{112}+O(h^9),\\
 \Logm(1+h)&=h-\frac{h^2}{2}+\frac{h^3}{3}-\frac{h^4}{4}+O(h^5),\\
 \sqrt{1+h}&=1+\frac h2-\frac{h^2}{8}+\frac{h^3}{16}
               -\frac{5h^4}{128}+O(h^5),\\
 \arccos(1-\delta)&=\sqrt{2\delta}
 \left(1+\frac\delta{12}+\frac{3\delta^2}{160}
                 +\frac{5\delta^3}{896}+\frac{35\delta^4}{18432}
                 +O(\delta^5)\right).
\end{align*}
For example, if $h=t^\lambda+t^\mu$ with $0<\lambda<\mu$, the first terms of $\sin h$ can be obtained by substituting this finite sum into its power series. Mixed terms such as $t^{2\lambda+\mu}$ must be collected with every other term of the same exponent. The displayed order in a polynomial truncation is not necessarily the order of all Hahn monomials: if $\mu\gg\lambda$, arbitrarily many powers $t^{n\lambda}$ precede $t^\mu$. The support convention handles this automatically and does not replace it by a single real small parameter.

\begingroup
\small
\setlength{\parskip}{0pt}
\begin{thebibliography}{99}
\setlength{\itemsep}{0.35em}
\raggedright
\addcontentsline{toc}{section}{References}

\bibitem{SourceAnalysis}
\emph{Surcomplex analysis}, supplied manuscript,
\texttt{surcomplex\_analysis(2).tex}, September 21, 2026.
Used for normal-form conventions, Hahn summability, finite Taylor lifting, coherent integrals, and the distinction between fine-local and coherent analyticity. Sections on global fine-analytic exponentials provide explicit predecessors of the infinite-phase examples.

\bibitem{SourceFinite}
\emph{Finite Geometry of Hahn-Coherent Surcomplex Zeros: Conservation of Multiplicity, Normal Forms, and Collision-Stable Residues}, supplied manuscript in
\texttt{surcomplex\_finite\_geometry.zip}, September 21, 2026.
Relevant background for positive-support conservation of finite zero clusters; its several-variable theorems are not prerequisites for the triangle geometry here.

\bibitem{SourceIntersections}
\emph{Finite Intersections in Surcomplex Analysis: Support-Controlled Division, Residue Duality, and Sharp Root Stability}, supplied manuscript in
\texttt{surcomplex\_intersections.zip}, September 21, 2026.
Relevant background for support-controlled division, exact local intersection algebras, and inverse-Jacobian valuation loss near colliding roots.

\bibitem{Conway}
John H. Conway, \emph{On Numbers and Games}, Academic Press, 1976;
2nd ed., A K Peters, 2001.

\bibitem{Gonshor}
Harry Gonshor, \emph{An Introduction to the Theory of Surreal Numbers},
London Mathematical Society Lecture Note Series \textbf{110},
Cambridge University Press, 1986.
\href{https://doi.org/10.1017/CBO9780511629143}{doi:10.1017/CBO9780511629143}.

\bibitem{Alling}
Norman L. Alling, \emph{Foundations of Analysis over Surreal Number Fields},
North-Holland Mathematics Studies \textbf{141}, North-Holland, 1987.

\bibitem{vdDE}
Lou van den Dries and Philip Ehrlich,
``Fields of surreal numbers and exponentiation,''
\emph{Fundamenta Mathematicae} \textbf{167} (2001), 173--188.
\href{https://doi.org/10.4064/fm167-2-3}{doi:10.4064/fm167-2-3}.
Erratum, \emph{Fundamenta Mathematicae} \textbf{168} (2001), 295--297.
The uses here are restricted analytic evaluation, real closedness, and the exponential group law, not the ordinal estimates addressed by the erratum.

\bibitem{Neumann}
B. H. Neumann, ``On ordered division rings,''
\emph{Transactions of the American Mathematical Society} \textbf{66} (1949), 202--252.

\bibitem{Poonen}
Bjorn Poonen, ``Maximally complete fields,''
\emph{L'Enseignement Math\'ematique} (2) \textbf{39} (1993), 87--106.
\href{https://math.mit.edu/~poonen/papers/amsval.pdf}{Author's full text}.
The algebraic-closure input is Corollary 4.

\bibitem{DLMF}
F. W. J. Olver et al., eds., \emph{NIST Digital Library of Mathematical Functions},
Chapter 4, especially \S\S4.14, 4.21, 4.23--4.24.
\href{https://dlmf.nist.gov/4}{dlmf.nist.gov/4}.
Consulted September 21, 2026. Reference for ordinary trigonometric identities and inverse-function expansions; the surreal domain and support arguments are supplied in this article.

\bibitem{DR}
Michael A. Dritschel and James Rovnyak, ``The operator Fej\'er--Riesz theorem,''
arXiv:0903.3639, 2009.
\href{https://arxiv.org/abs/0903.3639}{arXiv record}.
Cited for the established factorization problem and its classical context; the scalar real-closed-field proof used here is given in full.

\bibitem{CE}
Ovidiu Costin and Philip Ehrlich, ``Integration on the surreals,''
\emph{Advances in Mathematics} \textbf{452} (2024), article 109823.
\href{https://arxiv.org/abs/2208.14331}{arXiv:2208.14331}.
This is a separate integration framework, not an identification with the elementary coefficientwise functional of \cref{sec:arcs}.

\bibitem{BM}
Alessandro Berarducci and Vincenzo Mantova,
``Surreal numbers, derivations and transseries,''
\emph{Journal of the European Mathematical Society} \textbf{20} (2018), 339--390.
\href{https://doi.org/10.4171/JEMS/769}{doi:10.4171/JEMS/769}.
Cited to distinguish a derivation on the scalar field from the function derivatives used here.

\end{thebibliography}\endgroup
\end{document}
